\section{Exactness, amenability, and Ramsey properties}

\subsection{The group-theoretic exactness criterion}

\begin{theorem}[Proved: Lead 1]\label{A:coinvariants}
In setting (S), let $R$ be a unital ring. Write $\mathcal I_M$ for
the directed poset of nonempty finite strong subsets of $M$, with
inclusions. For $V:A_M\to R\text{-}\mathrm{Mod}$ put
$S_RV=\operatorname{colim}_{B\in\mathcal I_M}V(B)$. This is a smooth
$R[G]$-module under
\[
 g\,[B,v]=[gB,V(g|_B)v].
\]
The functor $S_R$ is exact, and naturally
\[
 \operatorname{colim}_{A_M}V\cong(S_RV)_G.
\]
Consequently $\operatorname{colim}_{A}$ is exact over $R$ if and only if
coinvariants preserve monomorphisms of smooth $R[G]$-modules. In
particular, exactness over a specified ring depends only on the
topological group $G$, not on its chosen strong Fra\"iss\'e presentation.
No rigidity, strong-amalgamation, or functorial-amalgamation hypothesis
is needed here.
\end{theorem}
\begin{proof}
Every finite subset of $M$ is contained in a finite strong stage of its
rich-chain construction. Thus $\mathcal I_M$ is directed. The displayed
action respects the filtered-colimit equivalence relation, because an
inclusion $B\subseteq C$ is carried by $g$ to $gB\subseteq gC$ and the
resulting square of embeddings commutes. Its group law follows from
functoriality. The class $[B,v]$ is fixed by $G_B$, so the action is
smooth. Exactness of filtered colimits of modules proves exactness of
$S_R$, including its asserted variance.

Each arrow $f:B\to C$ of $A_M$ extends, by strong homogeneity, to some
$g\in G$, and factors as $B\xrightarrow{g|_B}gB\hookrightarrow C$.
The relations defining the colimit over $A_M$ are therefore precisely
the inclusion relations followed by the coinvariant relations.
Explicitly, the mutually inverse maps send $[B,v]$ in the categorical
colimit to the coinvariant class of $[B,v]$ and conversely. This argument
also checks finite sums whose summands originally occur at different
objects: directedness first carries their finitely many supports to
one stage.

For a smooth $R[G]$-module $N$, define $T_N(B)=N^{G_B}$. If
$f:B\to C$ extends to $g$, set $T_N(f)n=gn$. Two extensions differ on
the right by an element of $G_B$, so this is well-defined. Moreover
$G_C\subseteq gG_Bg^{-1}$, which proves that the value lies in
$N^{G_C}$. The filtered colimit of $T_N$ is $N$: every vector is fixed
by some open subgroup, hence by $G_B$ for some finite strong $B$.
An injection $N'\hookrightarrow N$ induces objectwise injections
$T_{N'}\hookrightarrow T_N$. Exactness of $\operatorname{colim}_A$
therefore implies injectivity on smooth coinvariants. The reverse
implication follows by applying the exact functor $S_R$ and then
coinvariants to an objectwise short exact sequence. Coinvariants are
right exact, so preserving monomorphisms is exactly their exactness.
\end{proof}



\subsection{The convex and compact fixed-point rows}

\begin{theorem}[Proved: Lead 2, including the strong-class passage]
\label{A:convex}
In setting (S), the following are equivalent: amenability of $G$;
the finite-function convex Ramsey property stated in R12(ii), with
every embedding required to be strong; (Re); and (CC). Rigidity is
not required. The convex Ramsey statement is Moore's
Theorem~7.1(5)--(6) in the published paper, numbered
Theorem~6.1 in the supplied arXiv source, for ordinary ages; the proof below establishes the
strong version directly, rather than treating cofinality as a transfer
theorem \cite{ABMoore2013}.
\end{theorem}
\begin{proof}
We give the three comparisons. For amenability implying the infinite
convex Ramsey statement, fix strong $P,Q\subset M$ and finitely many
colourings $c_i:\operatorname{Emb}_{\leq}(P,M)\to[0,1]$. In the compact
product of these colouring spaces take the orbit closure of their
tuple under precomposition by $G$. This is a continuous $G$-flow:
each coordinate has an open stabilizer. An invariant probability
measure on that closure has a barycentre whose coordinates, for each
$i$, are constant, since $G$ acts transitively on strong embeddings
of $P$. Approximate this barycentre on the finitely many coordinates
$\operatorname{Emb}_{\leq}(P,Q)$ by a finite convex combination of
orbit points. Approximate those orbit points in turn by translates
of the original tuple. The resulting finite convex combination of
embeddings of $Q$ into $M$ has arbitrarily small oscillation for all
$c_i$. Its finite support is contained in a finite strong stage.

To obtain a single finite target working for every tuple of
colourings, suppose that no target works for a fixed positive
tolerance. The closed subsets of
$([0,1]^n)^{\operatorname{Emb}_{\leq}(P,M)}$ consisting of colourings
for which every probability measure supported in a given finite
strong target has oscillation at least that tolerance have the
finite intersection property. Here the minimum over a finite
probability simplex is continuous in the finitely many colour
coordinates. A global bad colouring contradicts the infinite
statement at any smaller tolerance. This proves the finite version.

Conversely let $G$ act continuously on a compact Hausdorff space $X$.
To produce an invariant probability measure it suffices, by
compactness of the probability-measure space, to make finitely many
invariance equations for finitely many continuous real functions
$f_i$ arbitrarily accurate. Fix the group elements $h_1,\ldots,h_s$
appearing in those equations. Uniform continuity of the action on
compact $X$ gives an open identity neighbourhood $U$ such that
$|f_i(ux)-f_i(x)|<\eta$ for all $x\in X$, all $u\in U$, and all $i$.
Choose a finite strong $P$ with $G_P\subseteq U$, and a finite strong
$Q$ containing $P\cup\bigcup_jh_jP$. Fix $x_0\in X$.

For each embedding $a:P\to M$, choose an extending automorphism
$g_a$ and colour $a$ by the rescaled values
$f_i(g_a^{-1}x_0)$. Any other extension differs on the right by
$G_P$, so these values differ by at most $\eta$ before rescaling.
Apply convex Ramsey to the embeddings $P\hookrightarrow Q$ and
$h_j|_P:P\to Q$. Extend the finitely many embeddings $b:Q\to M$
in the resulting probability measure to $g_b\in G$, and put the
same weights on $g_b^{-1}x_0$. The colour comparison is precisely
the difference between the averages of $f_i(x)$ and
$f_i(h_j^{-1}x)$, with an additional error at most $2\eta$.
Letting all tolerances tend to zero and using compactness yields an
invariant measure. This proves amenability.

For the comparison with (Re), enumerate the finite family
$a_0,\ldots,a_k:P\to Q$. At a finite target $R$, let $\Delta_R$
be the probability simplex on $A(Q,R)$ and put
$D\nu=(\nu a_i-\nu a_0)_{i=1}^k$. Finite-dimensional minimax gives
\[
 \min_{\nu\in\Delta_R}\sum_i\|\nu a_i-\nu a_0\|_1
 =\max_{\lambda_i\in[-1,1]^{A(P,R)}}
   \min_{\nu\in\Delta_R}\sum_i\langle\lambda_i,\nu a_i-\nu a_0\rangle.
\]
Convex Ramsey for $k$ functions makes the right side arbitrarily
small, after rescaling $[-1,1]$ to $[0,1]$. Notice that the number
of functions is $k$, independent of the size of the target; this
checks the potentially delicate quantifier interchange. Conversely,
(Re) controls every bounded function by the total-variation bound,
so implies convex Ramsey.

Finally we prove (Re)$\Leftrightarrow$(CC) directly for any nonempty
category with finite hom-sets and common successors. If (Re) fails
for $a_0,\ldots,a_k$, minimax gives an $\epsilon>0$ such that, for
every object $B$, the compact convex set
\[
 K(B)=\left\{(\lambda_i)\in[-1,1]^{kA(P,B)}:
 \sum_i\bigl(\lambda_i(ga_i)-\lambda_i(ga_0)\bigr)\geq\epsilon
 \text{ for all }g:Q\to B\right\}
\]
is nonempty. If there are no such $g$, the conditions are empty.
Restriction along $B\to C$ defines a continuous affine map
$K(C)\to K(B)$. A compatible point would make each $\lambda_i$
constant on all embeddings of $P$, by restriction along $P\to B$,
contradicting the inequality at $B=Q$, $g=1_Q$. Thus (CC) fails.

For the reverse direction, let $K:A^{\mathrm{op}}\to\mathrm{CompConv}$
have nonempty values. Fix finitely many compatibility equations and
choose a common successor $Q$ of their objects, with maps $j_B:B\to Q$.
For each source $B$ the comparisons are the finite family
$j_B,j_Cf:B\to Q$. Successive applications of (Re), followed by
convolution of the probability measures, make all these finitely
many families close simultaneously: subsequent pushforwards
contract total variation, so earlier bounds persist. Choose a
point $z\in K(R)$ at the final target and set
\[
 x_B=\sum_{g:Q\to R}\nu(g)K(gj_B)z.
\]
For any continuous seminorm on the ambient locally convex space of
$K(B)$, compactness bounds that seminorm uniformly on $K(B)$.
The total-variation estimates therefore make the chosen equations
arbitrarily accurate. Their closed-entourage versions have the
finite intersection property in $\prod_BK(B)$, and compactness
gives an exact compatible point. This proves (CC).
\end{proof}

\begin{theorem}[Known and rechecked: the compact Hausdorff row]
\label{A:compact-ramsey}
In setting (S) with (Rig),
\[
 \mathrm{(CH)}\quad\Longleftrightarrow\quad A\text{ is Ramsey}
 \quad\Longleftrightarrow\quad G\text{ is extremely amenable}.
\]
Each condition implies integral exactness. Finite hom-sets hold
because the structures are finite, even when the language is
countably infinite.
\end{theorem}
\begin{proof}
We rechecked T A.12--A.14: if finite-pattern Ramsey fails, the
nonempty compact sets of bad colourings form a contravariant
diagram with no compatible point. Conversely, simultaneous finite
Ramsey requirements solve any finite set of compatibility equations
to any compact uniform entourage; compactness solves them all.
With finite hom-sets a finite pattern can be the entire hom-set,
so finite-pattern and morphism Ramsey agree. Joint embedding gives
the common successors required there. The finite-family integral
affine-equalizer proof was rechecked in T A.6: a character of the
cokernel obstruction into $\mathbb R/\mathbb Z$, coloured into
$(3k+1)^k$ boxes, contradicts Ramsey unless the augmentation ideal
is all of $\mathbb Z$.

The last equivalence is exactly the strong-class formulation in
Evans--Hubi\v cka--Ne\v set\v ril, Theorem~2.13, which cites KPT
Theorem~4.8 and Nguyen Van Th\'e Theorem~1
\cite{ABEHN2019,ABKPT2005}. It uses only the extension of finite
strong partial isomorphisms and the cofinal basis $G_B$; these are
the same two properties used explicitly in the preceding proof.
Thus no transfer from a reduct's age is being presumed.
\end{proof}

\begin{proposition}[Proved: dense subgroups and compact averaging]
\label{A:averaging}
Let $G$ be any nonarchimedean topological group and $\Gamma\leq G$ a
dense subgroup. For every smooth module $N$, the quotient maps give a
natural equality $N_\Gamma=N_G$. If $G$ has an increasing sequence of
compact subgroups with dense union, smooth rational coinvariants are
exact. Thus Lead 5(a) holds in setting (S).
\end{proposition}
\begin{proof}
For $n\in N$ and $g\in G$, the open coset $gG_n$, where $G_n$ here
denotes the stabilizer of the vector $n$, meets $\Gamma$. Thus $gn$ is
equal to $\gamma n$ for some $\gamma\in\Gamma$; the two submodules of
coinvariant relations coincide.

If $K$ is compact, every orbit in a smooth discrete $K$-module is
finite. For any vector $v$, its stabilizer has finite index and contains
an open normal subgroup, obtained by intersecting finitely many
conjugates. Averaging over the resulting finite quotient defines
$e_Kv$; refinements of that quotient give the same value. This defines
a natural idempotent $e_K:N\to N^K$ over $\mathbb Q$, with kernel
spanned by $kv-v$. Consequently $N_K\cong N^K$ naturally, and averaging
proves exactness. For an increasing sequence $K_i$, the quotient
$N_{\bigcup_iK_i}$ is the filtered colimit of the $N_{K_i}$. Filtered
colimits preserve injections, and the first assertion identifies this
quotient with $N_G$.
\end{proof}

\begin{proposition}[Proved: the finite-index obstruction]
\label{A:finite-index}
In setting (S), if $V\leq U\leq G$ are open and $[U:V]=m<\infty$,
there is an injective morphism of smooth $R[G]$-modules
\[
 R[G/U]\longrightarrow R[G/V],\qquad
 [gU]\longmapsto\sum_{uV\in U/V}[guV].
\]
It induces multiplication by $m$ on coinvariants. In particular,
exactness over $\mathbb F_p$ requires $p\nmid m$. Exactness over all
prime fields implies (Fin). If a finite member of the class has an
automorphism of order $m>1$, it has no affine equalizer with the identity
over any ring in which $m$ is not a unit. Hence integral exactness,
quasi-filteredness, and totality each imply (Rig).
\end{proposition}
\begin{proof}
The coset sums have pairwise disjoint supports, so their map is
injective over every ring, including rings with zero divisors. Their
supports also show equivariance, and all basis vectors have open
stabilizers. Both transitive permutation modules have coinvariants
$R$, with each basis vector mapping to $1$. The induced map is therefore
$r\mapsto mr$; for $R=\mathbb F_p$ and $p\mid m$, this is the zero map
on a nonzero module.

Let $\sigma\in\operatorname{Aut}(X)$ have order $m$. Precomposition by
$\sigma$ acts freely on $A(X,Y)$: $f\sigma^j=f$ and injectivity of $f$
give $\sigma^j=1$. An element $t\in R[A(X,Y)]$ with $t\sigma=t$ has
constant coefficients on these free $m$-cycles, whence
$\operatorname{aug}(t)\in mR$. Augmentation one forces the central
integer $m$ to be a unit. For $R=\mathbb Z$ this is impossible, so IV
fails. The implication from totality to quasi-filteredness is the
all-action criterion audited in the background returns.
\end{proof}

\begin{proposition}[Proved: the discrete case of Lead 1]
\label{A:discrete}
For a discrete group $H$ and a nonzero unital commutative ring $R$,
coinvariants on all $R[H]$-modules are exact if and only if $H$ is
locally finite and the order of every finite subgroup is a unit of
$R$. For $R=\mathbb Q$ this is local finiteness. For $R=\mathbb Z$
it forces $H=1$.
\end{proposition}
\begin{proof}
Exactness passes to every subgroup $L\leq H$: induction
$R[H]\otimes_{R[L]}-$ is exact because $R[H]$ is a free right
$R[L]$-module, and its $H$-coinvariants are the $L$-coinvariants of the
original module. If $L$ is finitely generated, its augmentation ideal
is finitely generated, so the trivial module $R$ is finitely presented
over $R[L]$. Exactness of coinvariants says that this right module is
flat. A finitely presented flat module is projective, so the
augmentation has an $L$-invariant section $1\mapsto t\in R[L]$.
A nonzero invariant vector of finite support is constant on every
left-translation orbit, and the single orbit is $L$ itself. Thus $L$
is finite, and $1=\operatorname{aug}(t)=|L|r$ for its common coefficient
$r$. Conversely, each finite subgroup admits the averaging idempotent,
and coinvariants by a locally finite group are the filtered colimit
of these exact functors.
\end{proof}

\subsection{Amenable groups with a translational obstruction}

\begin{example}[Proved: Lead 6(a)]\label{A:translation}
Let $M=(\mathbb Z;(d_j)_{j\in\mathbb Z})$, where
$d_j(x,y)$ means $y-x=j$, and take all embeddings between finite
nonempty induced substructures. This is a countable rigid Fra\"iss\'e
class with finite hom-sets and automorphism group $\mathbb Z$ with the
discrete topology. Its group is amenable, but its category is not
exact over any nonzero ring, is not quasi-filtered, and is not total.
Strong amalgamation, (St), and (Op) all fail.
\end{example}
\begin{proof}
Every nonempty partial isomorphism is the restriction of a unique
translation. Joint embedding follows by taking the union of two
translated finite sets. For amalgamation over a nonempty base,
translate the two extensions so that their identified base agrees and
take their union. This proves amalgamation, although extra points can
be forced to coincide. Specifically, two copies of $\{0,1\}$ over
$\{0\}$ have their second points forced to agree, so (SA) fails.
Homogeneity and the extension property in $\mathbb Z$ follow from the
same translation description. A translation preserving a nonempty
finite set is zero, proving rigidity. The stabilizer of a point is
trivial, so the topology on the automorphism group is discrete.
F\o lner intervals prove amenability.

Let $a_0,a_1:\{0\}\rightrightarrows\{0,1\}$ select the two points.
For any target $Y$, identify an embedding of $\{0,1\}$ with its initial
integer. If $t\in R[A(\{0,1\},Y)]$ equalizes $a_0,a_1$, extending its
coefficient function by zero to $\mathbb Z$ gives a finitely supported
function invariant under translation by $1$. It is zero, so it cannot
have augmentation one. This proves the claimed failures over every
nonzero ring. For two disjoint finite nonempty supports, both pointwise
stabilizers are trivial whereas the stabilizer of their intersection
is $G$; thus (St) fails. A proper nonzero subgroup $m\mathbb Z$ is open
and is neither $G$ nor a nonempty finite-set pointwise stabilizer, so
(Op) fails.
\end{proof}

\begin{example}[Proved: Lead 6(b)]\label{A:fibred-translation}
Let $M=\mathbb Z\times\mathbb Q$, with the equivalence relation of
equal first coordinates, the rational order inside each class, and
binary relations $d_j$ specifying difference $j$ between first
coordinates. Its age with all embeddings satisfies (S), (Rig), (SA),
and (FA). Its automorphism group is topologically
\[
 \operatorname{Aut}(\mathbb Q,<)^{\mathbb Z}\rtimes\mathbb Z.
\]
It is amenable. The category has finite hom-sets and fails exactness
over every nonzero ring, quasi-filteredness, and totality; (St) and
(Op) fail.
\end{example}
\begin{proof}
An embedding translates all represented class coordinates by one
integer and is increasing within each represented class. Every finite
partial isomorphism extends by that translation and independent
automorphisms of the rational orders. The extension property is
equally immediate. Finite structures are rigid since a translation
preserving their finite class set is zero, and their finite orders
have no nontrivial automorphisms.

To amalgamate over a nonempty base $P$, first align class coordinates
using any point of $P$. In each common class amalgamate the two finite
orders over their common $P$-part, retaining distinct points off $P$.
Where this part is empty put the first side before the second; in
every gap of a nonempty part do the same. This yields a strong
amalgam, functorially for embeddings of the first side over $P$.
It proves (SA) and (FA), including classes represented on both sides
but absent from $P$. Joint embedding is obtained by choosing any
relative translation and applying the same construction.

The displayed topological semidirect product follows from the
explicit automorphisms; basic stabilizers constrain finitely many
fibre maps and require zero translation. The group
$K=\operatorname{Aut}(\mathbb Q,<)^{\mathbb Z}$ is extremely amenable:
finite products are, and their finitely supported union is dense.
The quotient is the amenable discrete group $\mathbb Z$; the standard
fixed-point proof for compact convex spaces proves amenability of the
extension. Indeed, the nonempty compact convex $K$-fixed subset is
$G$-invariant, and $\mathbb Z$ has a fixed point in it.

Take a two-point structure with its points in consecutive classes.
Pushing an affine equalizer of its two singleton embeddings to the
free module on class coordinates again gives a finitely supported
translation-invariant function on $\mathbb Z$, with the same
augmentation. It must vanish. The open normal subgroup $K$ is not
$G_F$: nonempty $F$ forces an actual rational point to be fixed,
whereas $K$ moves every point; $G_\varnothing=G$ is larger than $K$.
For two disjoint finite nonempty supports the generated stabilizers
are contained in $K$, so (St) fails as well.
\end{proof}

\subsection{Signed averaging over a compact Ramsey expansion}

\begin{theorem}[Proved: the signed averaging criterion]\label{A:signed-averaging}
Let $G$ be a nonarchimedean topological group, $H\leq G$ a closed
co-precompact subgroup, and $X=\widehat{G/H}$ the completion for the
right uniformity. Let $k$ be a commutative ring. Suppose that smooth
$k[H]$-coinvariants are exact and that there is a finitely additive,
$G$-invariant function
\[
 \mu:\operatorname{Clop}(X)\longrightarrow k,\qquad \mu(X)=1.
\]
No positivity or bounded-variation condition is imposed on $\mu$.
Then smooth $k[G]$-coinvariants are exact. The same conclusion holds
for modules over any unital ring $R$ if $\mu$ is integer-valued and
smooth $R[H]$-coinvariants are exact.

If $k$ is a field and smooth $k[G]$-coinvariants are exact, such a
normalized invariant $k$-valued function exists. Consequently, when
$H$ has exact smooth $k$-coinvariants, it is an exact criterion.
In particular, an amenable $G$ possessing a co-precompact extremely
amenable subgroup has exact smooth rational coinvariants.
\end{theorem}
\begin{proof}
For an open subgroup $U\leq G$, the finitely many double cosets
$U\backslash G/H$ form a finite clopen uniform partition of $G/H$.
These partitions define its right uniformity; their completions are
clopen partitions of $X$. Thus every function on $G/H$ constant on
the left $U$-orbits extends uniquely to a finite-valued locally
constant function on $X$. This also proves that $X$ is compact and
zero-dimensional, and that left translation by $G$ extends to a
continuous action on it.

Let $N$ be a smooth $k[G]$-module. For $n\in N$, define
\[
 f_n:G/H\longrightarrow N_H,\qquad f_n(gH)=[g^{-1}n]_H.
\]
Changing $g$ to $gh$ does not change this class. If $U$ fixes $n$,
then $f_n(ugH)=f_n(gH)$ for $u\in U$. Hence $f_n$ extends to a
finite-valued locally constant function on $X$, still denoted $f_n$.
Its integral is a finite algebraic sum:
\[
 s_N(n)=\int_X f_n\,d\mu
       =\sum_{C\in\mathcal P}\mu(C)f_n(C)\quad\in N_H,
\]
where $\mathcal P$ is any finite clopen partition on which $f_n$ is
constant. Finite additivity proves independence of the partition.
Common refinements show that $s_N:N\to N_H$ is $k$-linear and
natural in smooth $G$-module maps. For $a\in G$,
$f_{an}(gH)=f_n(a^{-1}gH)$, so invariance of $\mu$ gives
$s_N(an)=s_N(n)$. It therefore descends to a natural map
$\bar s_N:N_G\to N_H$.

Let $q_N:N_H\to N_G$ be the quotient. Its value on $f_n(gH)$ is
always $[n]_G$. Normalization gives
$q_N\bar s_N([n]_G)=[n]_G$. Thus the $G$-coinvariant functor is a
natural direct summand of the $H$-coinvariant functor on smooth
$G$-modules. In particular, for an injection $N'\hookrightarrow N$,
the injection $N'_H\hookrightarrow N_H$ and the two natural sections
show that $N'_G\to N_G$ is injective. This proves exactness. Integer
weights act centrally on every left $R$-module, so the same proof
proves the assertion for an arbitrary unital $R$.

For the converse let $C_{\mathrm{lc}}(X,k)$ be the smooth $G$-module
of locally constant $k$-valued functions, with $k$ discrete. Each
function has finite image, by compactness. Each of its finitely many
clopen fibres is stabilized by an open identity neighbourhood, by
continuity of the action and compactness of $X$; intersecting these
neighbourhoods gives an open stabilizer of the function. Exactness
therefore sends the injection of constant functions
$k\hookrightarrow C_{\mathrm{lc}}(X,k)$ to an injection on
coinvariants. Since $k$ is a field, choose a $k$-linear functional
on $C_{\mathrm{lc}}(X,k)_G$ taking the class of $1$ to $1$.
Its value on $[1_C]$ defines $\mu(C)$. The characteristic-function
identities give finite additivity, and the coinvariant quotient
gives invariance.

Finally, if $G$ is amenable, $X$ has an invariant Borel probability
measure. Choose a $\mathbb Q$-linear map
$\ell:\mathbb R\to\mathbb Q$ with $\ell(1)=1$, by extending
$\{1\}$ to a vector-space basis. Composing the measure of clopens
with $\ell$ gives a normalized invariant rational signed function.
Extreme amenability of $H$ gives exactness over $\mathbb Q$ by the
finite-pattern Ramsey criterion for a nonarchimedean presentation,
and the first part applies. The passage through $\ell$ does not
preserve positivity; positivity is not needed in this argument.
\end{proof}

\begin{lemma}[Proved: integral weights on finite linear extensions]
\label{A:poset-weights}
Let $(P,\prec,<)$ be a finite poset with an independent reference
linear order. For a linear extension
$\pi=(x_1,\ldots,x_n)$ of $\prec$, put
\[
 w_P(\pi)=\sum_{\mathcal B}(-1)^{n+|\mathcal B|},
\]
where $\mathcal B$ ranges over the partitions of this list into
nonempty consecutive blocks such that every block is a
$\prec$-antichain and its elements occur increasingly in $<$.
For $P=\varnothing$ the empty partition has weight $1$. These
integers are natural under isomorphisms preserving both relations,
have total sum $1$, and satisfy, for every induced $Q\subseteq P$
and every linear extension $\tau$ of $Q$,
\[
 \sum_{\pi|_Q=\tau}w_P(\pi)=w_Q(\tau).
\]
\end{lemma}
\begin{proof}
We prove the restriction identity when $P=Q\cup\{x\}$; iteration
proves the general case. A convenient bookkeeping device assigns a
real coordinate $t_y$ to each vertex and imposes
$t_y<t_z$ whenever $y\prec z$. Sorting by $(t_y,y)$, with ties
broken by the reference order, yields a linear extension. A block
partition in the formula is exactly a cell on which equal
coordinates are the blocks, with strictly increasing block heights.
A cell with $b$ blocks has assigned sign $(-1)^b$; the factor
$(-1)^{|P|}$ is applied after summing these signs. This description
is finite combinatorics: no measure or convergence is involved.

Fix a cell of $Q$ with $b$ blocks and block heights
$s_1<\cdots<s_b$. The admissible values of $t_x$ form the nonempty
open interval
\[
 \left(\max_{y\prec x}t_y,\ \min_{x\prec z}t_z\right),
\]
where a missing bound means $-\infty$ or $+\infty$. If both bounds
exist, transitivity gives $y\prec z$ for each such pair, so the
lower bound is strictly smaller than the upper bound. Suppose
exactly $m$ of the existing heights lie strictly inside this
interval. Splitting at those heights gives $m+1$ open-interval
possibilities, each a cell with $b+1$ blocks, and $m$ equality
possibilities, each a cell with $b$ blocks. The elements equal to
$x$ are necessarily incomparable with it, and their placement in
the block is uniquely determined by the reference order. Thus all
these cells are precisely the cells contributing above the fixed
cell of $Q$, counted once. Their signed sum is
\[
 (m+1)(-1)^{b+1}+m(-1)^b=(-1)^{b+1}.
\]
Multiplication by $(-1)^{|P|}=(-1)^{|Q|+1}$ gives the original
weight $(-1)^{|Q|+b}$ of the cell of $Q$. Summing over its cells
proves the restriction identity. Restriction to the empty set
proves normalization. Naturality is immediate from the formula.

Equivalently, the formula is the compactly supported Euler
characteristic of the indicated cells, multiplied by $(-1)^{|P|}$.
The elementary interval calculation above supplies the entire
additivity argument needed here. The attached exhaustive integer
check covers every reference-ordered labelled poset on at most
five vertices and all $102085$ one-vertex marginal equations;
it is an independent finite check, not the proof for arbitrary $P$.
\end{proof}

\begin{theorem}[Refuted: A.3(i), A.3(iv), and the first clause of A.3(v)]
\label{A:ordered-posets}
Let $\mathcal P$ be the class of all finite structures $(P,\prec,<)$,
where $\prec$ is a strict partial order and $<$ is an independent
linear order, with all induced embeddings as strong embeddings.
Its category $A_{\mathcal P}$ satisfies (S), (Rig), (SA), (FA),
(FA$\varnothing$), (St), and (Op). It is exact over every unital ring, and hence is
total. Its automorphism group is not amenable and is therefore not
extremely amenable. Consequently rational exactness does not imply
amenability, integral exactness with rigidity does not imply extreme
amenability, and (CT) does not imply (CC), even in this rigid
finite-relational Fra\"iss\'e setting.
\end{theorem}
\begin{proof}
Hereditary closure and countability are immediate. Given two finite
posets over an induced common base, take their union and close the
two partial orders transitively. Every path leaving and re-entering
one factor can be shortened across the common base, so this adds
no relation inside a factor and creates no directed cycle. The
images therefore form a strong amalgam. The linear orders are
amalgamated by putting the first-side points before the second-side
points in each gap of the base. For an embedding of the first
factor over the base, the same shortening argument shows that the
induced map of amalgams is an induced embedding of posets, and the
gap convention makes it increasing. The construction is functorial.
It also works over the empty base, where the partial order is the
disjoint union and the reference order is concatenation. This
proves all the named amalgamation hypotheses; rigidity follows
from the reference order. Proposition~\ref{B:poset-bridge} gives
the finite bridge proving (St), and Corollary~\ref{B:poset-op}
then proves (Op).

Expand by a further linear order $\triangleleft$ extending
$\prec$. This expansion class is Ramsey by Bodirsky,
Corollary~1.4(3), explicitly the class of finite posets with a
linear extension and an additional arbitrary linear order
\cite{ABBodirsky2014}. Its strong amalgamation is the same
construction, with the extension order also amalgamated in gaps.
Every chosen linear extension on a subposet extends to the whole
finite poset: repeatedly insert a new point between all its
predecessors and successors. Therefore the reduct of its
Fra\"iss\'e limit is the generic $(\prec,<)$-structure $M$.
Let $H$ be its expanded automorphism group. It is extremely
amenable and its smooth coinvariants are exact over every ring,
by the Ramsey integral criterion and extension of scalars.

The compact space $X$ of all linear extensions of $\prec$ on $M$
is the completion of $G/H$. Indeed its basic clopens specify the
extension on a finite subset. Expansion homogeneity identifies
two points in $G/H$ modulo the pointwise stabilizer of that subset
exactly when the specified finite extensions agree. There are
finitely many of them, all realized, and the orbit is dense by
the finite extension property. Thus $H$ is co-precompact.

For a finite $P\subseteq M$ and an extension $\pi$ of $P$, give
the cylinder $\{\triangleleft:\triangleleft|_P=\pi\}$ the integer
weight $w_P(\pi)$ of Lemma~\ref{A:poset-weights}. The restriction
identities make these prescriptions a well-defined finitely
additive function on all clopens: every clopen is a finite union
of cylinders at one finite stage, by compactness. Naturality makes
it $G$-invariant and normalization makes its total weight one.
Theorem~\ref{A:signed-averaging} proves exactness over every ring.
This is also a uniform existence proof for an affine equalizer of
every finite parallel family, by the Isbell--Mitchell criterion;
it does not merely solve a prescribed family.

Nonamenability is known for exactly this class from
Kechris--Soki\'c, Remark~3.2, pp.~78--79
\cite{ABKechrisSokic2012}. Here is a separate finite proof. Suppose
that $\lambda$ were a $G$-invariant probability measure on $X$.
Let $p$ be the probability that $x\triangleleft y$ for an
incomparable pair in reference order $x<y$. Homogeneity makes
this independent of the pair. Take a three-point reference order
$a<c<b$ with $a\prec b$ and $c$ incomparable with both. Since
every extension has $a\triangleleft b$,
\[
 \lambda(a\triangleleft c\triangleleft b)
 =\lambda(a\triangleleft c)-\lambda(b\triangleleft c)
 =p-(1-p)=2p-1.
\]
The same reference order with the reverse partial-order relation
$b\prec a$ gives $1-2p\geq0$. Thus $p=1/2$. It follows that,
for any $a\prec b$ and any $c$ incomparable with both, in whatever
reference order,
\[
 \lambda(a\triangleleft c\triangleleft b)
 =\lambda(a\triangleleft c)-\lambda(b\triangleleft c)=0.
\]
Fix a three-point antichain $A$. For each of its vertices $c$,
adjoin a fourth point above the other two in the partial order
and incomparable with $c$. Choose any reference position for this
fourth point. In that four-point poset, if $c$ is the middle of
the three antichain vertices in $\triangleleft$, then $c$ lies
between one of the lower vertices and their common upper bound.
The preceding equality makes this event have probability zero.
Restriction and homogeneity give the same conclusion on the
original antichain. Each of its three vertices thus has probability
zero of being the middle, although exactly one must be the middle.
This contradiction proves that $X$ has no invariant probability
measure, and hence that $G$ is not amenable. The attached program
also certifies infeasibility of the nonnegative marginal equations
through four vertices, with an exactly checked rational Farkas
certificate. The integral signed function above is allowed negative
values, so it does not contradict this argument.
Finally, exactness implies totality by TOTASPH C.6, whose
linearized-action argument we rechecked. The compact torsor
criterion in Part~C identifies (CT) with exactness, whereas
Theorem~\ref{A:convex} identifies (CC) with amenability.
These give the asserted refutations.
\end{proof}


\subsection{Rational exactness with no finite-index obstruction}

\begin{theorem}[Refuted: the converses in A.3(vi) and A.3(vii)]
\label{A:finite-cocycles}
Fix $r\geq2$ and a nonzero finite abelian group $\Lambda$.
The cocycle class $\mathcal A_r(\Lambda)$ of R13 satisfies (S),
(Rig), (SA), (FA), (St), (Op), (WEI), and (Fin). Its group is
amenable. For every commutative unital ring $R$, its colimit functor
is exact if and only if $|\Lambda|$ is a unit in $R$.
It is not total and does not satisfy (FA$\varnothing$).

In particular $\mathcal A_2(\mathbb F_p)$ is rationally exact and
satisfies (Fin), but fails exactness over $\mathbb F_p$ and over
$\mathbb Z$. Thus absence of finite-index open-subgroup obstructions,
even in this rigid setting with (SA), (FA), (St), and (Op), does not
make rational exactness imply prime-field or integral exactness.
\end{theorem}
\begin{proof}
The cocycle definitions, the finite extension argument, the
nonempty-base functorial amalgam, and the stabilizer calculation
are audited in Proposition~\ref{B:cocycle-audit} below. In
particular every finite subset is strong, every point outside
a finite $F$ has infinite $G_F$-orbit, and every open subgroup is
$G_F$. To check (Fin), write $V=G_C\leq U=G_B$ with $[U:V]<\infty$.
Each point of $C$ then has finite $G_B$-orbit, so $C\subseteq B$.
But this gives $G_B\leq G_C$, hence $U=V$. (WEI) is immediate
from (Op).

Expand a cocycle object $(S,c)$ by a primitive
$b\in C^{r-1}(\Delta^S;\Lambda)$ with $db=c$. The expansion class
is precisely the class of arbitrary $(r-1)$-cochains on finite
linearly ordered sets: the relation $c=db$ is determined by the
primitive. Encode a nonzero value of $b$ on an increasing
$r$-tuple by one of the finitely many relation symbols indexed by
$\Lambda\setminus\{0\}$, and use no such relation for value zero.
These are finite coloured $r$-uniform hypergraphs, with no
constraint except mutual exclusion of different colours on one
tuple. Their unordered class has free amalgamation: all new
cross-tuples have value zero. Adding an arbitrary linear order is
Ramsey by Evans--Hubi\v cka--Ne\v set\v ril, Theorem~1.3
\cite{ABEHNFunctions}. The expansion group $H$ therefore has exact
smooth coinvariants over every ring.

We check the reduct and the finite fibres explicitly. Every
cocycle on a simplex has a primitive. Given $S\subseteq T$ and a
primitive $b_S$ of $c_T|_S$, choose a primitive $b_T^0$ on $T$.
Then $b_S-b_T^0|_S$ is an $(r-1)$-cocycle. Since $r-1>0$, it is
$da_S$ for an $(r-2)$-cochain on $S$. Extend $a_S$ arbitrarily to
$T$ and put $b_T=b_T^0+da_T$. This extends $b_S$ and has
$db_T=c_T$. Thus the reduct of the expanded limit is the given
cocycle limit. Restriction between primitive sets is a surjective
map of finite torsors under the homomorphism of cocycle groups.
All its fibres have the same cardinality, a power product of the
prime divisors of $|\Lambda|$.

Let $X$ be the compact space of all global primitives of the limit
cocycle. Nonemptiness follows by the same extension argument and
compactness. As in the ordered-poset example, its finite cylinders
identify $X$ with $\widehat{G/H}$. Uniform measure on the finite
primitive set over each finite $S$ is compatible with restriction,
because the fibres just checked have equal sizes. These cylinder
measures extend to a $G$-invariant Borel probability measure on
$X$. They also give an invariant normalized
$\mathbb Z[1/|\Lambda|]$-valued function on its clopens. Theorem
\ref{A:signed-averaging}, followed by the ring homomorphism into
$R$, proves exactness whenever $|\Lambda|$ is invertible in $R$.
For amenability, let $Y$ be any compact $G$-flow. Extreme
amenability of $H$ gives an $H$-fixed $y$. The map
$gH\mapsto gy$ is right uniformly continuous: whenever $g'=ug h$
with $u$ near the identity and $h\in H$, the two values are $ugy$
and $gy$, uniformly close by compactness of $Y$. It extends
continuously to $X$. Pushing its invariant probability measure
to $Y$ proves amenability.

Here is a finite named obstruction proving the converse. Choose a
prime $p\mid|\Lambda|$ and a surjection $\chi:\Lambda\to\mathbb F_p$.
Choose $\lambda\in\Lambda$ with $\chi(\lambda)=1$. Let $Q$ be the
ordered set $\{0,\ldots,2r+1\}$. Define an $(r-1)$-cochain $b_Q$
to take value $\lambda$ on the $r$-tuple $(1,\ldots,r)$ and zero
on all other increasing $r$-tuples, and give $Q$ the cocycle
$c_Q=db_Q$. Its first block $Q_1=\{0,\ldots,r\}$ has cocycle
value $\lambda$; its second block
$Q_0=\{r+1,\ldots,2r+1\}$ has value zero.
Let $P$ be the ordered $r$-point structure, whose $r$-cocycle is
necessarily zero. For $\epsilon=0,1$ and $0\leq i\leq r$, let
$a_{\epsilon i}:P\to Q$ be the increasing embedding onto the
$i$th face of $Q_\epsilon$. All $2(r+1)$ maps are arrows.

Suppose at a target $D$ there were an affine equalizer
$t\in R[A(Q,D)]$ of this family, where $p$ is not a unit in $R$.
Pass to the nonzero $\mathbb F_p$-algebra $\bar R=R/pR$.
Choose a primitive $b_D$ of $c_D$ on $D$ and evaluate
$\chi b_D$ on the free $\bar R$-module on embeddings of $P$.
If all $t a_{\epsilon i}$ equal one vector $v$, the two alternating
sums
\[
 \sum_{i=0}^r(-1)^i(\chi b_D)(t a_{\epsilon i})
\]
have the same value, namely
$(\sum_i(-1)^i)(\chi b_D)(v)$. But the cocycle identity says that
the sum for $\epsilon=1$ is $\operatorname{aug}(t)=1$ and the
sum for $\epsilon=0$ is zero. This contradicts $1\ne0$ in
$\bar R$. Thus the displayed family has no affine equalizer at
any target, for any such ring. If $|\Lambda|$ is not a unit, one
of its prime factors is not a unit, completing the classification.

Finally $BA_r(\Lambda)\simeq K(\Lambda,r)$ has nonzero $H_r$,
as audited below. It is not contractible, hence not total.
Functorial joint embedding would contract the base by the usual
identity--amalgam--constant natural transformations, so
(FA$\varnothing$) fails. This last failure is precisely why this
example does not settle A.3(ii) or A.3(iii).
\end{proof}

\subsection{A correction to the independently ordered equivalence example}

\begin{proposition}[Refuted: the odd-characteristic clause of FE-1]
\label{A:equivalence-primes}
Let $A_{E,<}$ be the category of finite sets with an equivalence
relation and an independent linear order, with all embeddings. For
every prime $p$, this category is not exact over $\mathbb F_p$.
More generally, for every integer $m\geq2$ it has an explicit
$m$-arrow family without an affine equalizer over any ring in which
$m$ is not a unit. In particular, the assertion in A.4(3) that it is
exact over all odd prime fields is false.
\end{proposition}
\begin{proof}
Let $B$ have $m$ increasingly ordered pairwise inequivalent points.
For the specific two-arrow family in FE-1, the same argument below
applies directly: the configuration is
$x_1<y_1<x_2<y_2$, with the two $x$'s equivalent and the two
$y$'s equivalent. Recording the ordered pair of image classes
sends $t\alpha_1$ to a vector $v$ on ordered distinct class pairs
and sends $t\alpha_4$ to its reversal. An equalizer makes $v$
invariant under this free involution, so
$\aug t=\aug v\in2R$. Thus the original parity assertion is
correct, with exactly the stated four-point object and arrows.

We now give the obstruction in every characteristic.
Let $C$ consist of $m$ consecutive blocks of $m$ points. Index its
equivalence classes by $\mathbb Z/m\mathbb Z$, and in block $j$ list
their labels as $j,j+1,\ldots,j+m-1$. Let $\alpha_j:B\to C$ be
the increasing embedding onto block $j$.

For any finite target $D$, let $\mathscr O(D)$ be the set of ordered
$m$-tuples of distinct equivalence classes of $D$. Define
$\rho:R[A(B,D)]\to R[\mathscr O(D)]$ by recording the classes of the
image points; it preserves augmentation. If $t\in R[A(C,D)]$ and
$w=\rho(t\alpha_0)$, then $\rho(t\alpha_j)$ is the cyclic rotation
of $w$ by $j$. Consequently an equalizer makes $w$ invariant under
the free cyclic action of order $m$ on $\mathscr O(D)$. Coefficients
are constant on free orbits, so $\operatorname{aug}(t)=
\operatorname{aug}(w)\in mR$. Augmentation one is impossible when
$m$ is not a unit. The argument covers every target and every ring;
the accompanying program checks the finite cyclic certificates for
$m=2,3,5,7$.
\end{proof}

\begin{proposition}[Proved: rational exactness of the same example]
\label{A:equivalence-rational}
The group of the independently ordered equivalence structure is
amenable but not extremely amenable. Its category is exact over
$\mathbb Q$, but not over $\mathbb Z$ or any $\mathbb F_p$, and is
neither quasi-filtered nor total. It satisfies (Rig), (SA), (FA), and
(FA$\varnothing$), but not (St), (Op), or (Fin).
\end{proposition}
\begin{proof}
The generic structure is a countable dense order partitioned into
countably many dense classes. Let $J$ be its class set. The action
homomorphism $G\to\operatorname{Sym}(J)$ is surjective and open.
Surjectivity follows by a back-and-forth order isomorphism realizing
any prescribed permutation of the dense colours. A finite point
stabilizer maps onto the pointwise stabilizer of the finitely many
colours it meets, proving openness. Its kernel $K$ preserves every
colour separately.

Finite linearly ordered sets carrying named colours form a Ramsey
class. One may reduce any instance to its finitely many colours and
apply the finite ordered-unary-relational Ramsey theorem, or the
usual finite product Ramsey proof. Thus $K$ is extremely amenable
and its smooth coinvariants are integrally exact, by the strong KPT
correspondence and the integral Ramsey criterion. The infinite
symmetric group has a dense union of finite symmetric groups, so
its smooth rational coinvariants are exact. For a smooth $G$-module
$N$, $N_K$ is smooth for $G/K$: the image of an open stabilizer is
open, since the quotient map just proved is open. Moreover
$N_G=(N_K)_{G/K}$. The composite of the two exact rational functors
is exact. Both groups in the extension are amenable, proving
amenability of $G$.

The action on the compact space of orders of $J$ has no fixed point,
as the quotient is the full symmetric group. Hence $G$ is not
extremely amenable. The preceding proposition gives all
characteristic obstructions. In particular quasi-filteredness fails,
so totality fails. Rigidity comes from the linear order. A functorial
strong amalgam merges precisely the equivalence classes linked
through the common base, takes the transitive closure of that
identification, and orders side points first-side-before-second-side
in each gap. Classes absent from the base remain distinct. This
construction also works over the empty base.

The stabilizer of one equivalence class is open and is not a finite
pointwise stabilizer: a point stabilizer in that class fixes an
actual point, while finitely many points outside it do not prevent
it from moving. Thus (Op) fails. If $x\ne y$ lie in the same class,
$\langle G_x,G_y\rangle$ preserves that class and is proper in $G$,
although $\{x\}\cap\{y\}=\varnothing$, disproving (St). Finally the
stabilizer of an unordered set of $m$ classes has the finite quotient
$S_m$, so (Fin) fails. The finite-index proof also explains why no
odd characteristic is exceptional.
\end{proof}

\subsection{Amenability of all abelian cocycle examples}

\begin{theorem}[Proved: amenability does not require finite cocycle coefficients]
\label{A:cocycle-amenability}
Let $\Lambda$ be a countable abelian group and $r\geq2$.
The automorphism group of the strong Fra\"iss\'e limit of
$A_r(\Lambda)$ is amenable. Thus $A_2(\Q)$ satisfies
(S), (Rig), (SA), (FA), (St), (Op), (WEI), and (Fin), has an
amenable automorphism group, and is exact over every prime
field, but is not exact over $\Q$ or $\Z$. It is
quasi-filtered but is not total, and it does not satisfy
(FA$\varnothing$). It has no co-precompact extremely amenable
subgroup.
\end{theorem}
\begin{proof}
We prove the finite-function convex Ramsey property, and then
apply Theorem~\ref{A:convex}. All cochains below use the displayed
linear order to orient their simplices. Fix finite cocycle
structures $P,Q$, embeddings $a_0,\ldots,a_t:P\to Q$, and
$\varepsilon>0$. Choose a primitive
$b_0\in C^{r-1}(\Delta^Q;\Lambda)$ of the cocycle on $Q$.
The restrictions $b_i=a_i^*b_0$ are primitives of the same
cocycle on $P$. Since positive-dimensional cohomology of a
simplex vanishes, for each $i$ choose
$q_i\in C^{r-2}(\Delta^P;\Lambda)$ with
\[
 b_i-b_0|_{a_0(P)}=d q_i.
\]
Here the notation on the left means that both restrictions are
pulled back to the fixed ordered vertex set $P$; it does not
identify different subsets of $Q$.

For any finite set $T\subseteq\Lambda$ and any $\delta>0$,
there is a finitely supported rational probability $\nu$ on
$\Lambda$ with
\[
 \|\tau_u\nu-\nu\|_1<\delta\qquad(u\in T).
\]
For completeness, list generators of the subgroup generated by
$T$, express each $u\in T$ as the image of a vector in
$\Z^d$, and push forward the uniform probability on
$[-N,N]^d\cap\Z^d$. Translation changes a proportion tending
to zero of this box. Pushforward cannot increase the
$\ell^1$ norm, so sufficiently large $N$ proves the assertion.

Choose such a $\nu$ for all coordinates of the finitely many
$q_i$, with error smaller than $\delta/m$, where
$m=\max(1,\binom{|P|}{r-1})$. On each $(r-1)$-element subset
of $Q$ choose a coordinate independently with law $\nu$,
giving a random $(r-2)$-cochain $a$ on $Q$. Put
\[
 b=b_0+d a\in\operatorname{Prim}(Q),
\]
and let $\mu_Q$ be its finite-support probability distribution.
For each $a_i$, the law of $a_i^*a$ is the same product law
$\nu^{\otimes\binom{|P|}{r-1}}$. The elementary telescoping
bound for product probabilities shows that translating this
product by $q_i$ changes its $\ell^1$ norm by less than
$\delta$. Applying the coboundary map and the fixed translation
by $b_0|_{a_0(P)}$ cannot increase that norm. Consequently the
primitive marginals
\[
 \mu_i=(a_i^*)_*\mu_Q\quad\hbox{on }\operatorname{Prim}(P)
 \quad\hbox{satisfy}\quad
 \|\mu_i-\mu_0\|_1<\delta.
\]
The assertions are also valid when the relevant cochain group
has no coordinates, in which case its distribution is a point
mass.

We spell out why these approximate primitive marginals produce
\emph{positive} convex averages on embeddings. The expansion
whose objects are arbitrary $(r-1)$-cochains is Ramsey: for
any given instance only finitely many cochain labels occur,
and the class of ordered structures with those independent
labels is a free amalgamation class. Theorem 1.3 of
\cite{ABEHNFunctions} applies. Taking the coboundary gives its
reduct into $A_r(\Lambda)$, and the primitive extension
argument in Theorem~\ref{A:finite-cocycles} applies without
any finiteness assumption on $\Lambda$.

Jointly embed the finitely many expanded objects
$(Q,b)$ with $b\in\operatorname{supp}\mu_Q$ into a finite
expanded object $X^*$, using embeddings $j_b$. There are only
finitely many expanded $P$-types among their restrictions
$a_i^*b$. Iterated finite Ramsey gives an expanded target
$Y^*$ with the following property: for any fixed finite family
of functions $f_j:A(P,Y)\to[0,1]$, there is one embedding
$g:X^*\to Y^*$ on whose image the values of every $f_j$
oscillate by less than $\eta$ on embeddings of each of these
expanded $P$-types. To obtain this form, partition $[0,1]$
into finitely many intervals of length less than $\eta$,
colour by the interval vector, and apply Ramsey successively
to the finitely many types. The target is chosen before the
functions are specified.

Define the actual probability on reduct embeddings by
\[
 \beta=\sum_{b\in\operatorname{supp}\mu_Q}
             \mu_Q(b)[g j_b]\in\R[A(Q,Y)].
\]
For every $i,j$, the expectation of $f_j$ under $\beta a_i$
differs by at most $\eta$ from the sum of one value per
expanded $P$-type against $\mu_i$. Those values lie in
$[0,1]$. Thus
\[
 |(\beta a_i)(f_j)-(\beta a_0)(f_j)|
       \leq2\eta+\|\mu_i-\mu_0\|_1
       <2\eta+\delta.
\]
Choose $2\eta+\delta<\varepsilon$. This proves the strong
finite-function convex Ramsey property with positive
coefficients, and hence amenability.

For $\Lambda=\Q$, the structural assertions are proved in
Proposition~\ref{B:cocycle-audit}. The complete field profile
is Corollary~\ref{C4:field-profile}; in particular prime-field
exactness gives quasi-filteredness by
Proposition~\ref{C:arithmetic-pairs}. Its classifying space is
$K(\Q,2)$, so it is not total. Under (FA) and (Op), adding
(FA$\varnothing$) would imply totality by
Theorem~\ref{B:op-total}, so that condition fails. Finally,
an amenable group with a co-precompact extremely amenable
subgroup has exact rational coinvariants by
Theorem~\ref{A:signed-averaging}; the rational obstruction
therefore excludes such a subgroup here.
\end{proof}

\subsection{An amenable total category that is not integrally exact}

% Independently checked structural portion of the ordered-vector candidate.
% Amenability and the Ramsey expansion are not asserted in this note.
\begin{proposition}[Proved: the structural hypotheses for punctured ordered vector spaces]
\label{A:ordered-vector-structure}
Fix a finite field $k$. Let $\mathcal V_k$ consist of finite-dimensional
$k$-vector spaces with zero omitted from the universe, equipped with an
arbitrary independent linear order. Include the empty structure for the
zero vector space. In the relational language use one relation for each
linear equation on nonzero vectors; require the strong embeddings to
extend to injective linear maps. Then (S), (Rig), (SA), (FA),
(FA$\varnothing$), (St), (Op), (WEI), and (Fin) hold. Its category of
nonempty objects is total.
\end{proposition}
\begin{proof}
There are countably many finite types. Strong subobjects are exactly
punctured linear subspaces with the induced order, and compositions and
the intermediate-substructure axiom in (S) hold. Pushouts of vector
spaces provide amalgamation with the two factors meeting exactly in the
base. Orders on their union extend because the two orders agree on
the intersection, and the other vectors may be placed afterward.
The linear order makes every finite object rigid.

We give a functorial version of that construction. For fixed
$P\leq Z$, select a vector-space complement $Z_0$ with $Z=P\oplus Z_0$.
For every $X\geq P$, take $X\oplus Z_0$. The union of its embedded
$X$ and $Z$ is ordered by putting the $X$-side before the $Z$-side
in each gap of the order on $P\setminus\{0\}$. The remaining vectors
are the pairs $(x,z)$ with $x\in X\setminus P$ and $0\ne z\in Z_0$.
Put all of these after the union, ordered first by any fixed order on
$Z_0\setminus\{0\}$ and then by the given order on $X\setminus P$.
For an increasing linear embedding $u:X\to Y$ fixing $P$, the map
$u\oplus\mathrm{id}$ preserves all three parts and their orders:
$u^{-1}(P)=P$, so no mixed vector becomes an old vector. It is strong
and the construction respects composition. This proves (FA), including
the case $P=0$, which is (FA$\varnothing$) in the punctured language.

Here is the finite bridge needed for (St). Let $B,C\leq D$ be finite
subspaces, $F=B\cap C$, and let $d_0,d_1:D\to E$ be prescribed
increasing linear embeddings agreeing on $F$. Take an even number
$s\geq8$ of factors around a circle, one factor $E$ and $s-1$ copies
of $D$. Identify alternate boundary subspaces $B,C$, using $d_0,d_1$
at the two boundaries of $E$. Form the vector-space colimit $Z_0$:
the direct sum of all factors modulo the linear span of their overlap
relations. We check carefully that every factor embeds.

Index a boundary coefficient by an element of its overlap subspace.
A linear combination of overlap relations has, in each interior
$D$-factor, the difference of its two adjacent boundary coefficients.
If this combination is supported in a single factor $T$, then at every
other interior factor the coefficients agree as vectors of $D$, and
therefore belong to $B\cap C=F$. If $T=E$, all boundary coefficients
are the same element of $F$, and the two contributions in $E$ cancel
because $d_0,d_1$ agree on $F$. If $T$ is an interior factor, at least
one of the two chains between $T$ and $E$ has an interior factor, so
propagation along it produces a coefficient $f\in F$. If the other
chain has interior factors, it similarly produces $f'\in F$; the
vanishing coefficient at $E$ is $d_0(f)-d_1(f')$, so $f=f'$. If that
other chain has no interior factor, its coefficient $x$ initially
lies only in the appropriate boundary subspace, but the equation at
$E$ is $d_i(x)=d_{1-i}(f)=d_i(f)$, so again $x=f$. In either case
the two coefficients adjacent to $T$ coincide. The coefficient in
$T$ is then zero as well. Thus the
relation subspace meets every factor trivially, as required.

We also need the exact identifications between the old factor sets in
order to extend their linear orders. The same calculation with support
in two factors gives this description. Along either arc between those
factors, all intervening boundary coefficients are equal and lie in
$F$ as soon as the arc crosses an interior $D$-factor. Consequently two
vectors in different factors become equal precisely through the
prescribed boundary identifications, allowing the two endpoints of
$E$ to use the actual equality of their images in $E$. Equivalently,
the set union of the embedded factors is the strict set quotient used
in the alternating bridge: $F$-vectors have full visibility, and each
other vector has a proper visibility interval of length at most three.
For adjacent factors the long complementary arc contributes only an
$F$-vector and adds no identification; for factors at distance two the
short arc can only pass through $E$; at larger distance both arcs
contribute only $F$. This covers every pair of factors.

Embed the finite order on $E\setminus\{0\}$ into $\mathbb R$. On the
nonzero vectors of $D$ interpolate between the real coordinates of
$d_0$ and $d_1$. These increasing maps are constant on $F$. As in
Proposition~\ref{B:poset-bridge}, subdivide the interval finely and
alternate a map retaining the old $B$-coordinates with the next full
interpolated map; the intermediate map already has the new
$C$-coordinates. A sufficiently fine subdivision preserves every
strict order gap. The coordinates descend through exactly the set
identifications just proved, giving a linear order on the union of
the factor sets after arbitrary tie breaking. Put the remaining
nonzero vectors of $Z_0$ afterward in any order. All factors are now
increasing strong subspaces of a finite member of $\mathcal V_k$.
This is the required alternating bridge with prescribed endpoints.

If $M$ is the Fra\"isse limit and $g\in G_{B\cap C}$, apply that bridge
to the identity on a finite strong $D\supseteq B+C$ and to $g|_D$,
inside a finite strong $E\supseteq D+gD$. Embed the bridge over $E$
into $M$ and extend each copy of $D$ to an automorphism. Consecutive
automorphisms differ alternately by point stabilizers of $B$ and $C$.
Hence $g\in\langle G_B,G_C\rangle$, which proves (St).

Finite strong supports are closed under intersections. The
minimal-support argument therefore applies: for an open $U$, minimize
the dimension of $F$ with $G_F\leq U$. For every $u\in U$, (St) gives
$G_{F\cap uF}\leq U$, so minimality implies $uF=F$. Since $u$ preserves
the finite independent order on $F\setminus\{0\}$, it fixes that set
pointwise, and $U=G_F$. Thus (Op), and therefore (WEI), hold.

For clarity, (Fin) also follows without an implicit algebraic-closure
assumption. Strong amalgamation gives arbitrarily many independent
copies over a finite strong $F$ of a finite strong extension containing
any chosen point outside $F$. Thus every $G_F$-orbit outside $F$ is
infinite. If $G_C\leq G_F$ has finite index, each point of $C$ has a
finite $G_F$-orbit and must lie in $F$. Therefore $G_C\supseteq G_F$,
so they are equal. By (Op) this handles every pair of open subgroups.
Finally (FA), (FA$\varnothing$), and (Op) give totality by
Theorem~\ref{B:op-total}.
\end{proof}


\begin{theorem}[Proved: amenability of independently ordered finite-field vector spaces]
\label{A:ordered-vector-amenability}
For every finite field $k$, the automorphism group of the limit of
$\mathcal V_k$ in Proposition~\ref{A:ordered-vector-structure} is amenable.
Equivalently, its category satisfies (Re) and (CC).
\end{theorem}
\begin{proof}
We prove (Re) by a finite random-colouring argument. Fix ordered
punctured spaces $P,Q$, of vector dimensions $c,d$ and underlying
cardinalities $p,q$, and list their increasing linear embeddings
$a_0,\ldots,a_t:P\to Q$. If there is at most one such embedding there
is nothing to prove. Fix $\epsilon>0$. Forget the vector structure
temporarily. The category of finite linear orders satisfies (Re),
by its Ramsey property and Theorem~\ref{A:convex}. Consequently there
are an integer $m$ and a probability measure $\nu$ on increasing
injections from the $q$-point order to $[m]$ such that, writing
$\nu_i=\nu a_i$ for its marginals on $p$-point order embeddings,
\[
 \sum_{i=1}^t\|\nu_i-\nu_0\|_1<\epsilon/3.
\]
We may take $\nu$ rational, by approximation within this strict
inequality. Extend $\nu$ and its marginals by zero to all maps to
$[m]$.

Put $\ell=|k|$ and choose a large $n\geq d$. Let $V=k^n$, and colour
every nonzero vector independently and uniformly by $[m]$. Denote
the random colouring by $\chi$. Order $V\setminus\{0\}$ by colour,
breaking ties by an arbitrary fixed order. For every injective
linear map $g:Q\to V$ assign the nonnegative weight
\[
 w(g)=\nu(\chi\circ g).
\]
Only increasing maps with distinct colours on their images have
positive weight. Let $W=\sum_gw(g)$. If $e_r(n)$ is the number of
injective linear maps $k^r\to k^n$, then
\[
 e_r(n)=\prod_{j=0}^{r-1}(\ell^n-\ell^j),\qquad
 \mathbb E W=e_d(n)m^{-q}>0.
\]

For a linear embedding $u:P\to V$, let
$X_i(u)=\sum_{g a_i=u}w(g)$ be the unnormalised $i$th marginal.
There are exactly $N=e_d(n)/e_c(n)$ extensions of $u$ along each
$a_i$. Conditional on the colours on $uP$, every such extension has
$q-p$ distinct remaining vectors outside $uP$. Therefore
\[
 H_i(u):=\mathbb E[X_i(u)\mid\chi|_{uP}]
       =N m^{-(q-p)}\nu_i(\chi\circ u).
\]
This formula includes all nonincreasing and repeated-colour patterns,
where both sides are zero.

We bound the conditional variance without assuming independence of
all extensions. Two extensions $g,h$ of $u$ have independent weights,
conditional on $\chi|_{uP}$, whenever $gQ\cap hQ=uP$. For a fixed $g$
and a uniformly chosen extension $h$, each vector of $Q\setminus a_iP$
has image uniformly distributed on $V\setminus uP$. A union bound
therefore gives
\[
 \Pr(gQ\cap hQ\ne uP)
 \leq \eta_n:={ (\ell^d-\ell^c)^2\over\ell^n-\ell^c}
 \longrightarrow0.
\]
Here the zero vectors are included when counting the displayed
subspaces; their removal has no effect on the difference sets.
Since $0\leq w(g)\leq1$, the absolute value of every covariance is
at most one, and the covariance is zero for the other pairs. Hence
\[
 \operatorname{Var}(X_i(u)\mid\chi|_{uP})\leq N^2\eta_n,
 \qquad
 \mathbb E|X_i(u)-H_i(u)|\leq N\sqrt{\eta_n}.
\]
Summing over all $e_c(n)$ linear maps $u$ yields
\[
 \mathbb E\|X_i-H_i\|_1\leq e_d(n)\sqrt{\eta_n}.
\]
Also the colours of the $p$ distinct points of $uP$ are independent,
so
\[
 \mathbb E\|H_i-H_0\|_1
   =e_d(n)m^{-q}\|\nu_i-\nu_0\|_1.
\]
Both norms are taken on all linear embeddings $P\to V$; the actual
marginals are supported on the increasing ones, so this convention
does not enlarge their norms artificially.

Let $D=\sum_{i=1}^t\|X_i-X_0\|_1$. By the triangle inequality,
\[
 \mathbb E D\leq e_d(n)
 \left(2t\sqrt{\eta_n}
       +m^{-q}\sum_{i=1}^t\|\nu_i-\nu_0\|_1\right).
\]
Choose $n$ so large that $2t m^q\sqrt{\eta_n}<\epsilon/3$.
Then $\mathbb E(D-\epsilon W)<0$. Thus some colouring satisfies
$D<\epsilon W$. Necessarily $W>0$, because $W=0$ makes every weight
and every marginal zero. Normalising the weights gives a probability
measure on increasing strong embeddings $Q\to V$ for which the sum
of the required total-variation differences is less than $\epsilon$.
This proves (Re), and Theorem~\ref{A:convex} proves amenability.
\end{proof}

% Standalone addition for Part A; no amenability or Ramsey claim is made here.
\begin{proposition}[Proved: an ordered vector-space parity obstruction]
\label{A:ordered-vector-parity}
Let $\mathcal V^{<}_2$ have as objects the nonzero points of finite
positive-dimensional $\mathbb F_2$-vector spaces, with an arbitrary
linear order, and as arrows the restrictions of injective linear maps
which preserve that order. This category is not exact over
$\mathbb F_2$ or over $\mathbb Z$. In fact the five embeddings
specified below have no affine equalizer over $\mathbb F_2$ at any
target.
\end{proposition}
\begin{proof}
Let $P$ be the punctured two-dimensional vector space, whose three
points are denoted $x_1<x_2<x_3$. Label the nonzero vectors of
$\mathbb F_2^3$ by $1,\ldots,7$ in binary, so addition is bitwise
exclusive-or, and give the resulting object $Q$ the reference order
\[
 1<2<3<4<6<5<7.
\]
Define $a_j:P\to Q$ by their increasing image triples:
\[
\begin{array}{c|ccccc}
 j&1&2&3&4&5\\\hline
 (a_jx_1,a_jx_2,a_jx_3)&(1,2,3)&(1,6,7)&(2,5,7)&(3,4,7)&(3,6,5).
\end{array}
\]
Each triple is the set of nonzero vectors of a two-dimensional
subspace. Since every permutation of the three nonzero vectors of
$\mathbb F_2^2$ is linear, these are arrows of the specified category.

At any target $D$, choose an auxiliary ordered vector-space basis
and let $\prec$ be its natural binary order: compare coefficient
vectors at the largest index where they differ, with $0<1$.
This auxiliary order is unrelated to the given reference order of
$D$. For a nonzero vector $v$, write $h(v)$ for its largest nonzero
basis index and let $\epsilon_v$ be the linear functional extracting
the coefficient at that index. For distinct vectors $x,y$,
\[
 \mathbf1_{x\prec y}=\epsilon_{x+y}(y)\quad\text{in }\mathbb F_2.
\]
For nonzero distinct $u,v$, a second useful identity is
\[
 \mathbf1_{u+v\text{ lies between }u\text{ and }v\text{ in }\prec}
   =\epsilon_u(v)+\epsilon_v(u).
 \tag{*}
\]
Indeed, if $h(u)=h(v)$, the sum $u+v$ precedes both and the right
side is $1+1=0$. If $h(u)>h(v)$, then $v$ precedes both $u$ and
$u+v$, and the sum lies between them precisely when
$\epsilon_v(u)=1$; also $\epsilon_u(v)=0$. The remaining case is
symmetric. These arguments work in any finite dimension.

For any linear embedding $g:Q\to D$, abbreviate $g(i)$ to $i$ in
the following calculation. The identities
$1+2=5+6=3$ and $1+6=2+5=7$ imply
\[
\begin{split}
 &\mathbf1_{2\prec1}+\mathbf1_{1\prec6}
       +\mathbf1_{2\prec5}+\mathbf1_{5\prec6}\\
 &\qquad=\epsilon_3(1)+\epsilon_7(6)
       +\epsilon_7(5)+\epsilon_3(6)
  =\epsilon_3(7)+\epsilon_7(3).
\end{split}
\]
By (*) the last expression is the indicator that $4=3+7$ lies
between $3$ and $7$. Thus the sum of these five indicators is zero
for every $g$, regardless of the reference order of $D$.

Define five functions on $\mathcal V^{<}_2(P,D)$ by
\[
\begin{array}{ll}
 f_1(e)=\mathbf1_{e(x_2)\prec e(x_1)},&
 f_2(e)=f_3(e)=\mathbf1_{e(x_1)\prec e(x_2)},\\
 f_4(e)=\mathbf1_{e(x_2)\text{ between }e(x_1),e(x_3)},&
 f_5(e)=\mathbf1_{e(x_3)\prec e(x_2)}.
\end{array}
\]
The preceding calculation says $\sum_j f_j(ga_j)=0$ for every
$g:Q\to D$. On the other hand $\sum_j f_j(e)=1$ for every
$e:P\to D$. To check the latter identity, the two equal functions
$f_2,f_3$ cancel, and the median indicator is
$\mathbf1_{e(x_1)\prec e(x_2)}+
\mathbf1_{e(x_3)\prec e(x_2)}$ modulo two.

If an affine $t\in\mathbb F_2[\mathcal V^{<}_2(Q,D)]$ equalized
the five $a_j$, write their common composite as $m$. Extend the
$f_j$ linearly. Then
\[
 0=\sum_j f_j(ta_j)=\sum_j f_j(m)
    =\operatorname{aug}(m)=\operatorname{aug}(t)=1,
\]
a contradiction. The affine criterion gives nonexactness over
$\mathbb F_2$; integral exactness would imply exactness over
$\mathbb F_2$ and is therefore impossible as well.

The attached exact finite computation enumerates all $168$ natural
orders of the seven-point space and verifies the identity. It also
checks the associated marginal matrix and all $30$ isomorphism types
of reference orders. The coordinate argument above, rather than a
bounded search over targets, proves nonexistence at every target.
\end{proof}

\paragraph{Scope of this result.}
Combined with the adjacent structural and amenability theorems, this
proposition refutes A.3(ii) and the total-category implication in
A.3(v). It does not settle rational exactness of this category or
A.3(iii); no Ramsey expansion is being assumed.


\begin{corollary}[Refuted: A.3(ii) and the second clause of A.3(v)]
\label{A:ordered-vector-counterexample}
For $k=\mathbb F_2$, the category $\mathcal V_k$ is total and satisfies
(S), (Rig), (SA), (FA), (FA$\varnothing$), (St), (Op), (WEI), and (Fin).
Its group is amenable, but its colimit functor is not exact over
$\mathbb F_2$ or over $\mathbb Z$. Thus (CC) together with totality
does not imply (CT), even under all the structural hypotheses in
A.3(ii). Its group is not extremely amenable.
\end{corollary}
\begin{proof}
The structural assertions and amenability were just proved. The
five-arrow family in Proposition~\ref{A:ordered-vector-parity} has no
affine equalizer over $\mathbb F_2$ at any target. It therefore has
none over $\mathbb Z$ either. The finite-family criterion gives both
failures of exactness. By Part C, integral exactness is equivalent
to (CT), whereas (Re) is equivalent to (CC). Finally extreme
amenability would imply the Ramsey property and integral exactness,
contradicting this obstruction. Rational exactness is proved separately
in Theorem~\ref{A:ordered-vector-rational-exactness}. Thus this same
total example also refutes the converses in A.3(vi) for $p=2$ and
A.3(vii), even with all the structural hypotheses listed above.
\end{proof}

\subsection{The remaining named examples and the comparison table}

\begin{proposition}[Known and rechecked: the linear--cyclic example R10]
\label{A:cyclic-classification}
The category $C$ of finite nonempty sets with an independent
linear order and cyclic order is quasi-filtered, but not total
and not exact over any nonzero commutative ring. Its group is
neither amenable nor extremely amenable. It satisfies (Rig),
(SA), (FA), (St), (Op), and (Fin), but not (FA$\varnothing$).
Its classifying space is simply connected and
$H_2(BC;\Z)\cong\Z$.
\end{proposition}
\begin{proof}
The structural and group assertions except (Fin) were proved
in Proposition~\ref{B:cyclic-supports}; we supply the homology
and exactness checks used there. Let $P$ be the two-point
object and let $Q$ have linear order $1<2<3<4$ and cyclic
order $(1,3,2,4)$. Write $a_{ij}:P\to Q$ for its increasing
pairs. For any target $D$, cut its cyclic order and set
$\chi(u)=1$ if the larger of the two images of $u:P\to D$
occurs first in the resulting circular reading, and zero
otherwise. For every embedding $g:Q\to D$, the four
possible restricted cuts give
\[
 \chi(ga_{12})+\chi(ga_{23})-\chi(ga_{13})=1,
 \qquad
 \chi(ga_{12})+\chi(ga_{24})-\chi(ga_{14})=0.
\]
If $t$ equalizes the five arrows
$a_{12},a_{23},a_{13},a_{24},a_{14}$, applying the linear
extension of $\chi$ gives a common value equal both to
$\aug t$ and to zero. Thus the augmentation ideal of every
target kernel is zero, proving nonexactness over every
nonzero coefficient ring under discussion. The identities
are checked for every cut in \texttt{code/A\_B\_certificates.py}.
They also give a direct convex obstruction: for a probability
$t$, subtracting the two identities gives
\[
 1=\chi(ta_{23})-\chi(ta_{13})
       +\chi(ta_{14})-\chi(ta_{24}),
\]
so at least one of the two differences has absolute value
at least $1/2$. This alone disproves the convex Ramsey
property and amenability, independently of the circle-flow
proof in Part B.

Let $X_n$ be the set of cyclic orders on $\{0,\ldots,n\}$,
with deletion faces. The semisimplicial simplex category is
exactly $C$. Its nerve is the subdivision with face
occurrences retained, so $BC\simeq |X|$. There is one
vertex, one edge, two triangles, and six tetrahedra, and
the alternating cellular boundary maps are
\[
 d_1=0,\qquad d_2=(1,1),\qquad
 d_3=0:\Z^6\longrightarrow\Z^2.
\]
The last formula follows by deleting each vertex of each
of the six cyclic orders and counting the two triangle
orientations with alternating signs; the same certificate
script prints all six zero columns. Either triangle kills
the generator of the one-skeleton, so $\pi_1(BC)=0$.
Also $H_2(BC;\Z)=\ker(1,1)\cong\Z$, which excludes
totality. By (FA) and (Op), every action component is
contractible or equivalent to $BC$, so all are simply
connected. The condition-I/IV action characterization,
rechecked in DF A.2 and used in Corollary~\ref{B:wei}, gives
quasi-filteredness. Thus this example does supply an
additional strict quasi-filtered versus total separation
once the stabilizers have been proved.

Finally, if $G_C\leq G_B$ has finite index, the finite
$G_B$-orbit of every point of $C$ is contained in $B$,
because strong duplicate amalgams make every orbit outside
$B$ infinite. Hence $C\subseteq B$ and the inclusion is
equality. By (Op) this proves (Fin) for all open pairs.
(FA$\varnothing$) would now force totality by
Theorem~\ref{B:op-total}, so it fails.
\end{proof}

\begin{proposition}[Proved: finite graphs with all induced embeddings]
\label{A:graphs}
The category of nonempty finite graphs with all induced
embeddings is exact over $\Q$, but not over $\Z$ or any
$\F_p$. It is neither quasi-filtered nor total. Its random
graph automorphism group is amenable but not extremely
amenable; (St) holds, while (Rig), (Op), and (Fin) fail.
It satisfies (SA), (FA), and (FA$\varnothing$).
\end{proposition}
\begin{proof}
The free graph amalgam, with no new cross edges, proves the
three amalgamation assertions and the usual Fra\"iss\'e
properties. Expand by a completely arbitrary linear order.
The expansion is Ramsey by Theorem 1.3 of
\cite{ABEHNFunctions}. Its rich limit has an extremely
amenable automorphism subgroup $H\leq G$, and its orbit
completion is the compact space $X$ of all linear orders
on the random graph: every finite order pattern on an
induced finite graph embeds in the expanded rich limit.
There are $|S|!$ orders on a finite set $S$. Giving each
weight $1/|S|!$ gives restriction-compatible rational
probabilities, since every order on a subset has the same
number of extensions. These define an invariant probability
on $X$ and a rational-valued function on its clopens.
The signed averaging theorem proves rational exactness.
For any compact $G$-flow, an $H$-fixed point and the orbit
map from $G/H$ push this invariant probability to that
flow, exactly as in Theorem~\ref{A:finite-cocycles}.
This proves amenability directly.

The edgeless graph on $p$ vertices has a cyclic automorphism
of order $p$. The free-orbit coefficient argument of
Proposition~\ref{A:finite-index} proves nonexactness over
$\F_p$, as well as failure of rigidity, quasi-filteredness,
and totality. There is no invariant linear order on the
random graph: every two-point induced graph has an
automorphism interchanging its vertices, which extends by
homogeneity. Thus the compact order flow has no fixed
point, excluding extreme amenability.

For completeness the free-amalgamation proof of (St) can
be checked directly. Given finite $A,B$ and
$g\in G_{A\cap B}$, choose successively
$h_1\in G_A$ and $h_2\in G_B$ so that $h_1g(A)$ is
freely placed from $B$ over $A$, and $h_2h_1g(A)$ is
freely placed from $A$ over $B$. The latter copy meets
$B$ in $A\cap B$ and is freely placed from $A$ over
$A\cap B$. Choose also $h_3\in G_B$ with $h_3(A)$
freely placed from $A$ over $B$. The two enumerated
copies have the same type over $A$, so some $h_4\in G_A$
maps the former to $h_3(A)$ with matching enumeration.
Then $h_3^{-1}h_4h_2h_1g\in G_A$, proving the
generation identity. This is Lemma 2.7(i) of
\cite{ABMacphersonTent}; its entire finite argument was
checked, and all free placements here are allowed graph
amalgams.

Choose distinct vertices $a,b$. Their setwise stabilizer
$U$ has the pointwise stabilizer $G_{\{a,b\}}$ as a
normal subgroup of index two, by homogeneity. Also
$\Fix(U)=\varnothing$: it is contained in
$\Fix(G_{\{a,b\}})=\{a,b\}$ and a transposition in
$U$ moves both. The subgroup $U$ is proper. Consequently
it cannot be a finite pointwise stabilizer, since such a
stabilizer fixes precisely its finite support. This proves
failure of both (Op) and (Fin).
\end{proof}

\begin{corollary}[Proved: the complete table requested in A.4]
\label{A:example-table}
The entries in the following table hold for the named
categories. The column $\F_2$ records exactness over that
field. In every row with a negative $\F_2$ entry except
where explicitly qualified elsewhere, the proofs above
actually give failure in every prime characteristic.
\begin{center}
\small
\begin{tabular}{lcccccccc}
\toprule
Class & Total & $\Z$ & $\Q$ & $\F_2$ & Amen. & Extr. & St & Op\\
\midrule
$\Cdelta$ & yes & no & no & no & no & no & yes & yes\\
Convex equivalence and order & yes & yes & yes & yes & yes & yes & no & no\\
Independent equivalence and order & no & no & yes & no & yes & no & no & no\\
Independent poset and order & yes & yes & yes & yes & no & no & yes & yes\\
Independent cyclic and linear order & no & no & no & no & no & no & yes & yes\\
Finite graphs & no & no & yes & no & yes & no & yes & no\\
$C(2,2,1)$ & yes & no & no & no & no & no & yes & yes\\
$C(3,3,1)$ & yes & no & no & no & no & no & yes & yes\\
Translation structure on $\Z$ & no & no & no & no & yes & no & no & no\\
Fibred translation structure & no & no & no & no & yes & no & no & no\\
\bottomrule
\end{tabular}
\end{center}
None of these categories is filtered.
\end{corollary}
\begin{proof}
The two equivalence rows follow from Theorem~\ref{B:equivalence-types},
the modular obstruction in Proposition~\ref{A:equivalence-primes},
and the rational averaging in Proposition~\ref{A:equivalence-rational}.
The poset row is Theorem~\ref{A:ordered-posets}, with its
bridge in Part B. The cyclic and graph rows are the two
preceding propositions. The translation rows are
Examples~\ref{A:translation}--\ref{A:fibred-translation}.
The three predimension rows follow from the full parameter
classification in Part D. Each category has two different
embeddings from a singleton into some object, and every
arrow is injective. They therefore cannot be equalized by
postcomposition, which proves failure of filteredness.
\end{proof}


\begin{corollary}[Proved: amenability with metrizable universal minimal flow]
\label{A:metrizable-umf-exact}
In setting (S), suppose that $G$ is amenable and its universal minimal
flow is metrizable. Then $A$ is exact over $\Q$. Consequently any
counterexample to A.3(iii) must have nonmetrizable universal minimal
flow. This conclusion does not require (Rig), (FA), (FA$\varnothing$),
or (Op).
\end{corollary}
\begin{proof}
The group $G$ is a closed subgroup of the permutation group of a
countable set. It has a canonical relational presentation: put a
relation on each orbit of each finite tuple. Its age is an ordinary
Fra\"iss\'e class and its automorphism group is still $G$.
Zucker's Theorem~8.14 in \cite{ABZucker2016} gives a reasonable
precompact Ramsey expansion with the expansion property when the
universal minimal flow is metrizable. The automorphism group $H$ of
its limit is a closed co-precompact subgroup of $G$, and is extremely
amenable by the embedding version of the Ramsey correspondence.
The positive invariant probability on the compact completion of
$G/H$ supplied by amenability, followed by a $\Q$-linear functional
$\R\to\Q$ taking $1$ to $1$, is a normalized invariant finitely
additive $\Q$-valued measure on its clopens. The signed averaging
argument of Theorem~\ref{A:signed-averaging} makes smooth
$G$-coinvariants exact over $\Q$. Proposition~\ref{A:coinvariants}
then applies to the original strong presentation of $G$.
\end{proof}


\begin{theorem}[Proved: totality transfers to a presentation with finite closures]
\label{A:totality-finite-closure}
Let $A$ and $A'$ be two presentations satisfying (S), with limits
$M,N$, and identify their automorphism groups by a topological
isomorphism $G$. Suppose $A$ is total. Suppose also that every
nonempty finite subset of $N$ has a least finite strong superset,
written $\operatorname{cl}_N(S)$. Then $A'$ is total.
In particular, totality is a group invariant among presentations
whose nonempty intersections of finite strong subsets are strong.
No rigidity, functorial amalgamation, or intersection assumption
is imposed on the source presentation.
\end{theorem}
\begin{proof}
We work with actual finite strong subsets in the two limits.
Proposition~\ref{B:no-open-quotients} implies (Fin) for $G$.
Let $D$ have objects $(B,C)$ with $B\le M$, $C\le N$ nonempty
finite and $G_B\le G_C$. A morphism is a pair of restrictions
of one $g\in G$, whose images lie in the respective targets.
It is determined by its restriction to $B$: two extensions of
that restriction differ by an element of $G_B$, which fixes $C$.
All the restrictions in question are strong. The category $D$
has common successors. Indeed, first choose finite strong sets
containing the unions in $M$ and $N$, and then enlarge the former
so that its pointwise stabilizer fixes the latter.

We first prove that $D$ is total. Enumerate all pairs
$(B_i,n_i)$ with $B_i\le M$ finite nonempty and
$n_i\in\operatorname{Fix}_N(G_{B_i})$. Every point of $N$
occurs, since its stabilizer is open. Put
\[
 S_k(B)=\{g n_i: i\le k,\ gB_i\subseteq B\},
 \qquad C_k(B)=\operatorname{cl}_N(S_k(B))
\]
when $S_k(B)\ne\varnothing$. The set $S_k(B)$ is finite:
there are only finitely many embeddings $B_i\to B$, and the
image of $n_i$ depends only on the restriction to $B_i$.
These sets are functorial for strong embeddings of $B$.
The least-closure property makes $C_k(B)$ finite and invariant
under $G_B$. Its permutation action has an open kernel of
finite index in $G_B$, so (Fin) makes it pointwise fixed.
Consequently
\[
 j_k:B\longmapsto(B,C_k(B))
\]
is a functor from the nonempty full cosieve
$A_k=\{B:S_k(B)\ne\varnothing\}$ into $D$.
Each $A_k$ is connected by common successors and total:
extend any of its set actions by the empty set outside this
cosieve and apply totality of $A$.

Let $L:D\to\Set$ be an action. A finite simplicial diagram
in $N\El_D(L)$ involves only finitely many objects $(B,C)$.
Choose $k$ so large that every pair $(B,c)$ with $c\in C$
appears in the enumeration. Then $C\subseteq C_k(B)$ at
every one of those objects, and the simultaneous identity
restrictions give compatible arrows
$(B,C)\to j_k(B)$. They give a nerve homotopy from the finite
diagram into $\El_{A_k}(L j_k)$. The latter has weakly
contractible components. Every sphere in the nerve can be
represented, after simplicial approximation, on such a finite
diagram. It therefore contracts in its component of
$\El_D(L)$. This proves totality of $D$.

Finally the projection $p:D\to A'$, $(B,C)\mapsto C$, is
homotopy final. For fixed $C\in A'$, an object of $C/p$ is
$(B,C',f:C\to C')$. Choose $h_f\in G$ extending $f$.
It determines $h_f^{-1}B$ with stabilizer contained in $G_C$.
An arrow represented by $g$ goes to the restriction of
$h_{f'}^{-1}g h_f$, which fixes $C$. This defines a retraction
of $C/p$ onto $E_{G_C}$ for the source presentation: the
inclusion sends $B$ to $(B,C,\id_C)$, and the maps induced
by $h_f$ give a natural transformation from inclusion followed
by retraction to the identity. The category $E_{G_C}$ is the
connected action category for $G/G_C$ on $A$, hence is weakly
contractible by totality of $A$.

For any $L':A'\to\Set$, the induced functor on categories
of elements has the same comma categories $C/p$: the label
at its target is forced to be $L'(f)x$. Thus it too is
homotopy final. Totality of $D$ proves totality of $A'$.
If nonempty intersections of finite strong subsets of $N$
are strong, intersect all finite strong supersets of a fixed
nonempty finite $S$ inside one such superset. Only finitely
many intersections are needed, and their result is its least
finite strong closure. This verifies the final assertion.
\end{proof}

\begin{theorem}[Proved: totality and full cofinal subcategories]
\label{A:cofinal-totality}
Let $A$ be a nonempty connected small category satisfying condition I.
Let $C\subseteq A$ be a full subcategory such that every object of $A$
has a morphism to an object of $C$. Then $A$ is total if and only if
$C$ is total. Mere objectwise cofinality is sufficient here; homotopy
finality is a conclusion, not an additional assumption.
\end{theorem}
\begin{proof}
First suppose $C$ is total. For $a\in A$, the category $a/C$ is
nonempty. Condition I extends any two arrows $a\to c,a\to c'$ to
a commuting square in $A$. Extend its target once more into $C$.
Thus $a/C$ is connected. It is the category of elements of the
$C$-action $A(a,-)|_C$, so totality of $C$ makes its nerve weakly
contractible. Consequently $C\hookrightarrow A$ is homotopy final.
For every $L:A\to\Set$ the induced functor
$\El_C(L|_C)\to\El_A(L)$ has the same comma categories $a/C$:
the target label of an arrow out of $(a,x)$ is forced by $x$.
The homotopy-finality criterion therefore proves totality of $A$.

Conversely suppose $A$ is total. It is quasi-filtered. Both
condition I and the integral affine-equalizer condition IV descend
to $C$: construct the required target in $A$ and postcompose all
its arrows into an object of $C$. Fullness puts every resulting
arrow between objects of $C$ in $C$. Common successors descend in
the same way. Hence $C$ is connected and quasi-filtered. The
all-action characterization rechecked in DF A.2 says that every
component of every $\El_C(L)$ is simply connected. It remains to
prove its integral acyclicity.

Fix $L:C\to\Set$. For $a\in A$, let $\mathcal F(a)$ be the
category with objects $(c,x,f)$, where $c\in C$, $x\in L(c)$,
and $f:c\to a$. A morphism to $(c',x',f')$ is an arrow
$h:c\to c'$ satisfying $f'h=f$ and $L(h)x=x'$.
Postcomposition in $A$ makes $\mathcal F$ a covariant diagram
of categories. For $d\in C$, its components are indexed by
$y\in L(d)$, and the component indexed by $y$ has terminal object
$(d,y,\id_d)$. Here fullness ensures that every arrow $c\to d$
is an arrow of $C$, so that this component description is valid.
It follows that
\[
 H_q(B\mathcal F(d);\Z)=0\quad(q>0),\qquad
 H_0(B\mathcal F(d);\Z)=\Z[L(d)].
\]
For each $a$ choose an arrow $a\to d$ with $d\in C$.
For $q>0$, the resulting map on $H_q(B\mathcal F(-);\Z)$ is
zero. These positive-homology diagrams are therefore locally
zero: every element is killed after some arrow. By
Proposition~\ref{G:moduleinjectivization}, all their derived
colimits over the total category $A$ vanish, including degree
zero. Meanwhile
$H_0(B\mathcal F(-);\Z)=\Z[\pi_0\mathcal F(-)]$
is free on a set action, so totality gives vanishing of its
positive derived colimits.

For completeness, the spectral comparison used here has a direct
bar proof. Form the first-quadrant bicomplex
\[
 K_{p,q}=\bigoplus_{a_0\to\cdots\to a_p}
                  \Z[N_q\mathcal F(a_0)],
\]
with the ordinary bar differential horizontally and nerve
differential vertically. Taking vertical homology and then
horizontal homology gives
\[
 E^2_{p,q}=L_p\colim_A H_q(B\mathcal F(-);\Z).
\]
For fixed $q$, a basis element of $K_{p,q}$ consists of a
$q$-simplex
$(c_0,x_0)\to\cdots\to(c_q,x_q)$ of $\El_C(L)$ followed by
$c_q\to a_0\to\cdots\to a_p$. Thus its horizontal complex
is a direct sum, over those $q$-simplices, of the nerve chain
complexes of $c_q/A$. Each has an initial identity object.
Taking horizontal homology first therefore leaves precisely
$\Z[N_q\El_C(L)]$ in degree zero and nothing in higher degrees.
The total complex computes $H_*(B\El_C(L);\Z)$. The preceding
vanishing on its other spectral sequence now gives
$H_n(B\El_C(L);\Z)=0$ for every $n>0$.
Each component is simply connected, so the first-nonzero-group
form of the Hurewicz theorem makes it weakly contractible.
This proves totality of $C$.
\end{proof}

\begin{corollary}[Proved: unrestricted group invariance of totality]
\label{A:totality-group-invariant}
In setting (S), totality of the strong embedding category depends
only on the topological isomorphism type of $G=\Aut(M)$.
Together with Theorem~\ref{A:coinvariants}, this answers both
questions in A.1 affirmatively. Neither assertion needs (Rig),
(SA), (FA), (FA$\varnothing$), (St), (Op), or an intersection
hypothesis.
\end{corollary}
\begin{proof}
Let $A,A'$ be two strong presentations with topologically
isomorphic automorphism groups, and suppose $A$ is total.
Write $N$ for the limit of $A'$. Give $N$ its canonical
relational structure by naming each $\Aut(N)$-orbit on every
finite tuple. Its age $\widehat A'$ consists of all nonempty
finite subsets with the induced canonical structure.
Every embedding between these finite subsets extends to an
automorphism: apply the orbit relation to a tuple enumerating
the entire domain. Hence the canonical structure is homogeneous,
its age satisfies (S), and its topological automorphism group
is exactly $\Aut(N)$. Its strong closure is the identity on
finite subsets. Theorem~\ref{A:totality-finite-closure} therefore
implies that $\widehat A'$ is total.

Use actual finite subsets to identify $A'$ with the full
subcategory of $\widehat A'$ on the finite strong subsets of
the original structure. Fullness holds because any canonical
embedding extends to an automorphism of $N$; the strong-image
and restriction axioms in (S) make its image strong inside the
target. Conversely every strong embedding extends by strong
homogeneity and hence is canonical. Every finite subset of
$N$ lies in a finite strong stage, so this full inclusion is
objectwise cofinal. The canonical age has condition I by its
amalgamation property. Theorem~\ref{A:cofinal-totality} now
implies totality of $A'$. Interchanging the presentations
proves the equivalence.
\end{proof}

\begin{corollary}[Proved: totality passes to dense continuous images]
\label{A:totality-dense-image}
Let $G,H$ be automorphism groups of presentations satisfying (S).
If $\rho:G\to H$ is a continuous homomorphism with dense image
and a presentation of $G$ is total, then every presentation of $H$
is total. In particular, none of the automorphism groups of the
cocycle limits for $\mathcal A_r(\Lambda)$, $r\ge2$, $\Lambda\ne0$, can be
a continuous quotient, or a dense continuous image, of a group
with a total presentation.
\end{corollary}
\begin{proof}
Let $M$ be the source limit and $N$ a target limit. First give $N$
its canonical relational presentation, so every finite subset is
strong. Regard $N$ as a smooth $G$-set through $\rho$. Density
implies that every embedding between finite canonical subsets
of $N$ is induced by an element of $G$: the set of elements of
$H$ agreeing with a given finite embedding is a nonempty open
set. Continuity says that every point stabilizer in this
$G$-action is open.

These two facts are exactly the target-action facts needed in
the proof of Theorem~\ref{A:totality-finite-closure}; faithfulness
of the target action is not used. More explicitly, use pairs
$(B,C)$ with $B\le M$ nonempty finite, $C\subset N$ nonempty
finite, and $G_B$ fixing $C$ through $\rho$. Enumerate the pairs
$(B_i,n_i)$ with $G_{B_i}$ fixing $n_i$, and use the same finite
sets $S_k(B)$. Here $C_k(B)=S_k(B)$ because the target closure
is the identity. Their nonempty-support cosieves and the identity
arrows $(B,C)\to(B,S_k(B))$ contract every finite sphere in
every action category on the pair category, exactly as before.

For the projection to the canonical target age, the comma at
$C$ retracts onto $E_U$ of the source, where
$U=\rho^{-1}(H_C)$. This is open. Given a target embedding
$f:C\to C'$, choose $g_f\in G$ which induces it. Transport
$B$ by $g_f^{-1}$, and transport an arrow represented by $g$
by $g_{f'}^{-1}gg_f$. The latter fixes $C$ under $\rho$ and
therefore belongs to $U$. These are the same retraction and
natural transformation as in Theorem~\ref{A:totality-finite-closure}.
Totality of the source contracts $E_U$, so the projection is
homotopy final on all action categories. The canonical target
age is total. Theorem~\ref{A:cofinal-totality} passes this to
the original strong target age, and
Corollary~\ref{A:totality-group-invariant} to every presentation
of $H$. The cocycle conclusion follows from their nonzero
$H_r(B\mathcal A_r(\Lambda);\Z)=\Lambda$.
\end{proof}


% Standalone addition, following the ordered-vector parity proposition.
\begin{proposition}[Proved: an integral free-valued diagram with torsion colimit]
\label{A:ordered-vector-integral-torsion}
There is a diagram $M:\mathcal V^{<}_2\to\Ab$ whose values are
finite-rank free abelian groups and whose colimit contains an element
of order exactly two. This assertion remains valid whether the empty
punctured zero-dimensional vector space is included or omitted.
\end{proposition}
\begin{proof}
For an object $B$, let $\mathcal T(B)$ be the finite set of natural
binary orders on its underlying vector space, with zero omitted.
An embedding $a:B\to C$ induces the restriction map
$r_a:\mathcal T(C)\to\mathcal T(B)$. Restricting a natural order to
a subspace again gives a natural order: an adapted echelon basis
orders the subspace by its successive leading coordinates. Define
\[
 N(B)=\mathbb Z^{\mathcal T(B)},\qquad
 N(a)(f)=f\circ r_a,\qquad
 M(B)=N(B)/\mathbb Z\mathbf1_B.
\]
These are covariant functors. Each $N(B)$ is free of finite rank,
and its constant function $\mathbf1_B$ is primitive, so $M(B)$ is
also free of finite rank.

The uniform probabilities on $\mathcal T(B)$ are restriction
consistent. Indeed every natural order on a subspace extends by
extending an adapted ordered basis. The linear automorphisms of the
ambient space stabilizing that subspace act transitively on its
natural orders, and restriction is equivariant for this action.
Consequently all restriction fibers have the same cardinality.
It follows that
\[
 \mu_B(f)=\frac1{|\mathcal T(B)|}
                \sum_{\tau\in\mathcal T(B)}f(\tau)
\]
defines a natural homomorphism $N\to\mathbb Q$, where the latter
is the constant diagram. Since $\mathcal V^{<}_2$ is connected,
it induces
\[
 \mu:\operatorname{colim}N\longrightarrow\mathbb Q,
 \qquad \mu([\mathbf1])=1.
\]
In particular the map $\mathbb Z\to\operatorname{colim}N$ induced
by constant functions is injective. Right exactness of colimits gives
\[
 \operatorname{colim}M
   =\bigl(\operatorname{colim}N\bigr)/\mathbb Z[\mathbf1].
 \tag{IT.1}
\]

Take the objects $P,Q$ and the five embeddings $a_j:P\to Q$ from
Proposition~\ref{A:ordered-vector-parity}. Regard its five indicator
functions $f_j$ as integer-valued functions of the auxiliary natural
order on $P$. The parity identities proved there say precisely that
\[
 \sum_{j=1}^{5}f_j=\mathbf1_P+2h_P,
 \qquad
 \sum_{j=1}^{5}N(a_j)(f_j)=2h_Q
 \tag{IT.2}
\]
for integer-valued functions $h_P\in N(P)$ and $h_Q\in N(Q)$.
Explicitly $h_P=(\sum_jf_j-1)/2$ and
$h_Q=(\sum_jN(a_j)(f_j))/2$; their integrality follows from those
identities at every natural order, rather than from an averaging
argument. Passing to the colimit in (IT.2) yields
\[
 [\mathbf1]=2v,
 \qquad v=[h_Q]-[h_P]\in\operatorname{colim}N.
\]
Applying $\mu$ gives $\mu(v)=1/2$. Thus $v$ cannot be an integer
multiple of $[\mathbf1]$. Its image in (IT.1) is nonzero and is
killed by two, so has order exactly two, as claimed. The attached
finite certificate verifies (IT.2) on the six natural orders of $P$
and the $168$ natural orders of $Q$, and gives
$\mu_P(h_P)=2/3$ and $\mu_Q(h_Q)=7/6$.
\end{proof}

\paragraph{Scope.}
This is torsion in the ordinary colimit of a pointwise free integral
diagram. It does not assert that $T_1(M)$ has torsion, or determine
any positive derived colimit of this diagram. Although the maps of
$N$ have nonnegative entries in the bases of point indicators,
passing to $M=N/\mathbb Z\mathbf1$ does not supply compatible
nonnegative matrices or lifts to homomorphisms of free monoids.


\begin{proposition}[Proved: open subgroups and dense images]
\label{A:open-subgroups}
Let $G$ be the automorphism group of a presentation satisfying
\textup{(S)}. If that presentation is total, every \textup{(S)}
presentation of every open subgroup $H\leq G$ is total.
For any fixed unital coefficient ring $R$, exactness likewise
passes from $G$ to an open subgroup, and from $G$ along a
continuous homomorphism with dense image to another group
having an \textup{(S)} presentation.
Consequently no open subgroup of a group with a total
presentation can have a nontrivial cocycle group
$\operatorname{Aut}(M_{\mathcal A_r(\Lambda)})$, $r\geq2$,
$\Lambda\ne0$, as a continuous dense image.
\end{proposition}
\begin{proof}
For totality, first use Corollary~\ref{A:totality-group-invariant}
to replace the presentation by the canonical relational age of
$G$ on its original countable set $M$. Its category $A$ has all
nonempty finite subsets of $M$ as objects and the embeddings
induced by $G$ as arrows. Define the action
\[
 L(B)=(G/H)^{G_B}.
\]
On an embedding induced by $g$, the action is left translation
by $g$; two choices differ on the right by $G_B$, so this is
well-defined. Put $C=\El_A(L)$.

This category is connected. Given $(B,gH)$ and $(D,hH)$,
choose $u\in G$ with $ugH=hH$, and a finite set containing
$uB\cup D$. The two induced arrows give a common successor
with label $hH$. It is also nonempty because $H$ is open.
Moreover $C$ is total. Indeed an arbitrary action $K:C\to\Set$
induces the action
\[
 \widetilde K(B)=\coprod_{\ell\in L(B)}K(B,\ell)
\]
on $A$, and its category of elements is exactly $\El_C(K)$.
Totality of $A$ therefore contracts every connected component
of the latter.

Consider the full subcategory $D$ of $C$ whose label is $H$.
Its objects are the nonempty finite sets $B$ with $G_B\leq H$;
its arrows are precisely the embeddings induced by $H$.
Every $(B,gH)$ is isomorphic to $(g^{-1}B,H)$, so $D\simeq C$.
Give $M$ its canonical relational structure for $H$. In this
structure the admitted sets $B$ with $G_B\leq H$, and all their
isomorphic copies, form a strong class satisfying \textup{(S)}.
To check the details, choose a finite $F$ with $G_F\leq H$.
Every finite set is contained in the admitted set $B\cup F$,
so the admitted sets are cofinal. They are upward closed and
invariant under $H$. Every embedding of canonical finite
$H$-structures extends to an element of $H$, by the relation
naming the orbit of an enumeration of its entire domain.
The strong restriction axiom is automatic for these induced
embeddings between admitted objects. Joint embeddings and
amalgams are obtained by aligning the base by such an element
of $H$ and taking the finite union inside $M$. Countability
is inherited from $M$. The limit is the canonical $H$-structure
on $M$, whose automorphism group is $H$ because $H$ is closed.
Thus $D$ is an \textup{(S)} presentation of $H$, and is total.
Corollary~\ref{A:totality-group-invariant} proves the assertion
for every presentation of $H$.

For exactness use Theorem~\ref{A:coinvariants}. If $H$ is
open, smooth induction
\[
 N\longmapsto R[G]\otimes_{R[H]}N
\]
is exact: $R[G]$ is free as a right $R[H]$-module on a choice
of coset representatives. It is smooth because each tensor
has finite support and its stabilizer contains a finite
intersection of conjugates of open stabilizers in $H$.
There is a natural identification
\[
 \bigl(R[G]\otimes_{R[H]}N\bigr)_G\cong N_H.
\]
Exactness of $G$-coinvariants therefore implies exactness of
$H$-coinvariants.

If $\rho:G\to H$ is continuous with dense image, restriction
of a smooth $H$-module to $G$ is smooth and exact. The two
groups have the same orbit on each vector: the set of elements
of $H$ carrying one vector to a specified translate is a
nonempty open coset, so meets $\rho(G)$. Consequently their
coinvariant quotients agree, proving dense-image descent of
exactness. The final totality assertion follows by combining
the open-subgroup result with
Corollary~\ref{A:totality-dense-image}.
\end{proof}

The open-subgroup qualification is necessary. The group
$\operatorname{Aut}(\mathbb Q,<)$ has an exact total Ramsey
presentation, whereas its closed subgroup of integer
translations is isomorphic to the discrete group $\mathbb Z$,
whose rational coinvariants are not exact by
Proposition~\ref{A:translation}. The translation subgroup is
closed because the image of $0$ determines its integer
translation parameter in the pointwise convergence topology.

\subsection{A quantitative form of the remaining rational question}

\begin{proposition}[Proved: the necessary failure of a linear filling bound]
\label{A:distortion-gate}
Assume (S), and suppose that $G$ is amenable but $A$ is not
exact over $\Q$. There is a finite parallel family
$a_0,\ldots,a_k:p\to q$, with $k\geq1$, for which the following
holds. Put $X_b=\operatorname{Emb}_{\leq}(b,M)$ and define the
smooth $G$-module homomorphism
\[
 D:\Q[X_q]\longrightarrow \Q[X_p]^k,
 \qquad [g]\longmapsto([ga_i]-[ga_0])_{i=1}^k.
 \tag{DG.1}
\]
Let $J=\operatorname{im}D$, give the target its sum of the
coefficient $\ell^1$ norms, and define on $J$ the quotient
seminorm
\[
 \|w\|_D=\inf\{\|t\|_1:t\in\Q[X_q],\ Dt=w\}.
 \tag{DG.2}
\]
Then $J$ is generated by one $G$-orbit. There is a well-defined
$G$-invariant $\Q$-linear functional $\lambda:J\to\Q$ with
$\lambda(Dt)=\operatorname{aug}t$, and a sequence $w_n\in J$ such
that
\[
 \lambda(w_n)=1,\qquad \|w_n\|_D=1,
 \qquad \|w_n\|_1\longrightarrow0.
 \tag{DG.3}
\]
In particular no constant $C$ satisfies
$\|w\|_D\leq C\|w\|_1$ on this submodule. Conversely, if every
finite-family image (DG.1) has such a linear filling bound,
amenability implies rational exactness.
\end{proposition}
\begin{proof}
Choose a family with no rational affine equalizer, using the
finite-family criterion. Every finite collection of embeddings
into $M$ has its images in one finite strong subset. Therefore
an element of $\ker D$ with nonzero augmentation could be moved
to one actual object of $A$ and divided by its augmentation to
give an affine equalizer there. It follows that
$\operatorname{aug}(\ker D)=0$. This proves that $\lambda$ is
well defined and linear. Equivariance of $D$ and invariance of
augmentation give its $G$-invariance. Homogeneity makes $X_q$
a transitive $G$-set, so the images of its basis vectors form
one generating orbit for $J$.

For every $w=Dt$ one has
$|\lambda(w)|=|\operatorname{aug}t|\leq\|t\|_1$;
taking the infimum gives $|\lambda(w)|\leq\|w\|_D$.
By the amenability--Reiter equivalence there are finitely
supported probability measures $\mu_n$ on embeddings of $q$
into finite targets, regarded as elements of $\Q[X_q]$, with
$\|D\mu_n\|_1\to0$. Rational coefficients may be used:
approximate the finitely many probability weights by rational
ones of the same total mass, and use
$\|D(t)\|_1\leq2k\|t\|_1$ to control the error.
For $w_n=D\mu_n$, its displayed preimage has norm one, while
$\lambda(w_n)=1$ forces every preimage to have norm at least
one. This proves (DG.3), and the last assertion follows by
contraposition.
\end{proof}

This is a precise remaining analytic obstruction, not an
automatic-bound theorem for oligomorphic groups. Boundedness
in a quotient or filling seminorm does not imply boundedness
in the ambient coefficient norm. In the locally compact
setting, \cite[Propositions~4.2 and 4.4]{ABArora2021} explicitly
distinguishes these issues and requires a projective
direct-summand hypothesis for its positive undistortion result.
Those hypotheses are not assumed here. No rational exactness
conclusion for the independently ordered vector-space example
follows from this proposition alone.

\subsection{Rational exactness of independently ordered binary vector spaces}

The rational entry for the ordered-vector example can be decided without
assuming a Ramsey theorem for two orders on vector spaces with arbitrary
linear embeddings.  We use a different category in which the embeddings
preserve a distinguished basis, and then transfer finite affine equalizers.

\begin{lemma}[Proved: Ramsey property with a distinguished ordered basis]
\label{A:ordered-basis-ramsey}
Let $\mathcal C$ have as objects finite positive-dimensional
$\mathbf F_2$-vector spaces $W$, a distinguished linearly ordered basis
$E_W$, and an independent linear order on $W\setminus\{0\}$.
Its arrows are injective linear maps taking basis elements to basis
elements and preserving both orders.  Then $\mathcal C$ is Ramsey,
has finite hom-sets and common successors, and is exact over every ring.
\end{lemma}
\begin{proof}
An object is equivalently a finite linearly ordered set $E$ with an
independent linear order on its nonempty subsets: a subset represents
the sum of its basis elements.  Fix objects $A,B$. If $A$ does not
embed into $B$, the Ramsey assertion is vacuous; otherwise
$|E_A|\leq |E_B|$. Put $b=|E_B|$.
For a nonempty subset $S=\{s_1<\cdots<s_k\}$ of size at most $b$, encode
$S$ by the $b$-tuple
\[
 \tau_b(S)=(s_1,\ldots,s_k,
                 \underbrace{s_k,\ldots,s_k}_{b-k\text{ entries}}).
\]
The distinct entries recover $S$, so this encoding is injective and
commutes with increasing injections.  On $E_A$ and $E_B$ put the
$b$-ary relation consisting of all such tuples.  Order that relation by
the given order on nonempty subsets, and retain the basis order as the
vertex order.

Hubi\v cka--Ne\v set\v ril's theorem on totally ordered structures
\cite[Theorem~4.46, arXiv v4]{DHubickaNesetril2019} gives a finite
Ramsey target $C_0$ for these two encodings.  Here ``totally ordered''
means that the vertices and the tuples of the relation carry independent
linear orders, both preserved by embeddings.  The quoted theorem applies
to one relation of arbitrary finite arity, including repeated entries.

In $C_0$, delete relation tuples that are not of the canonical form
$\tau_b(S)$, and add every missing canonical tuple.  Retain the order
on the remaining old relation tuples and extend it to the added tuples
arbitrarily.  This operation preserves every encoded copy of $B$:
on its vertex set all canonical tuples were already present and no
noncanonical tuple was present, because its map into $C_0$ was an
embedding.  Its two orders are unchanged.  Moreover every encoded
$A$-copy inside such a $B$-copy was already an $A$-copy in $C_0$.
Thus a coloring after completion restricts to the old $A$-copies, and
a monochromatic $B$ supplied by $C_0$ remains monochromatic afterward.

Finally give the nonempty subsets of the target basis of size greater
than $b$ arbitrary positions extending the order already defined on
the smaller subsets.  The resulting full powerset order gives an object
of $\mathcal C$ with the required Ramsey property.
There are finitely many basis injections between two objects.
Common successors are obtained by taking disjoint bases, merging each
of the two old orders, and placing the new mixed subset sums arbitrarily.
The Ramsey exactness theorem therefore gives exactness over every ring.
\end{proof}

\begin{lemma}[Proved: finite averaging through an exact category]
\label{A:exact-cover-averaging}
Let $u:\mathcal C\to\mathcal A$ be a functor, where $\mathcal C$ is
nonempty, has common successors, and is exact over a commutative ring $R$.
For $Y\in\mathcal A$ define
\[
 X_Y(C)=R[\mathcal A(Y,uC)].
\]
For a finite family $a_0,\ldots,a_k:P\to Q$, suppose there is
$\mu\in\operatorname{colim}_{\mathcal C}X_Q$ of augmentation one
whose images under every precomposition map
$a_i^*:\operatorname{colim}X_Q\to\operatorname{colim}X_P$ agree.
Then the family has an affine equalizer over $R$ in $\mathcal A$.
\end{lemma}
\begin{proof}
Let $D:X_Q\to X_P^k$ have components $a_i^*-a_0^*$.
Exactness of the colimit gives a surjection
\[
 \operatorname{colim}_{\mathcal C}\ker D
       \longrightarrow
 \ker\bigl(\operatorname{colim}_{\mathcal C}D\bigr).
\]
Lift $\mu$ to the left side.  A colimit element uses only finitely
many objectwise representatives.  Send their objects to one common
successor $C$ and add their images.  This yields
$t\in\ker D(C)\subseteq R[\mathcal A(Q,uC)]$ representing $\mu$.
Augmentation is natural, so $\operatorname{aug}t=1$.
The equality $D(C)t=0$ is precisely the desired affine equality.
\end{proof}

\begin{theorem}[Proved: rational exactness of the ordered-vector example]
\label{A:ordered-vector-rational-exactness}
Let $\mathcal A$ be the category of finite nonzero
$\mathbf F_2$-vector spaces with an arbitrary linear order on their
nonzero vectors and all order-preserving linear embeddings.
Then $\mathcal A$ is exact over $\Q$.
Consequently the example of
Theorem~\ref{A:ordered-vector-counterexample} is total, amenable, and
rationally exact, while it is not exact over $\mathbf F_2$ or over $\Z$.
\end{theorem}
\begin{proof}
Use the category $\mathcal C$ of
Lemma~\ref{A:ordered-basis-ramsey}, and let $u$ forget its distinguished
ordered basis.  First identify the components of the category of
elements of $C\mapsto\mathcal A(P,uC)$.
For an embedding $h:P\to uC$, let $U_h$ be the span of precisely those
distinguished basis vectors which occur with coefficient one in some
vector of $h(P)$.  Give $U_h$ the two induced orders.  Its isomorphism
type together with the marked map $h:P\to U_h$ is unchanged by arrows
of that category of elements.  Conversely $U_h\to C$ itself is such
an arrow.  Two marked closures of the same type therefore lie in the
same component, by their isomorphism and the two inclusions.  Thus
\[
 \operatorname{colim}_{\mathcal C}X_P
   =\Q[\{\text{isomorphism types of marked ordered basis closures of }P\}].
 \tag{BV.1}
\]
This uses only that the free-module functor preserves colimits.

Fix an arbitrary finite family $a_i:P\to Q$, and put
$p=\dim P$, $q=\dim Q$.  Let $Q^\vee$ be the binary linear dual and
form the binary simplex-code embedding
\[
 W=\mathbf F_2^{\,Q^\vee\setminus\{0\}},\qquad
 j:Q\longrightarrow W,\qquad
 j(v)=(\ell(v))_{\ell\in Q^\vee\setminus\{0\}}.
 \tag{BV.2}
\]
Its distinguished basis is the set of coordinate vectors, indexed by
the nonzero functionals $\ell$.  The map $j$ is injective because
linear functionals separate vectors.

Choose the basis order uniformly from all its possible orders.
Independently, order all vectors of
$W\setminus j(Q)$ uniformly, put that entire block before
$j(Q)\setminus\{0\}$, and put the latter block in the prescribed
order of $Q$.  Here and below the zero vector is omitted from all
orders.  Every such pair of choices gives an object $W_\omega$ of
$\mathcal C$ and an embedding $j:Q\to uW_\omega$.
There are finitely many choices; their uniform rational average defines
\[
 \mu=\frac{1}{|\Omega|}\sum_{\omega\in\Omega}[j:Q\to uW_\omega]
       \in\operatorname{colim}_{\mathcal C}X_Q,
 \qquad\operatorname{aug}\mu=1.
 \tag{BV.3}
\]

The basis closure $U_i$ of $ja_i(P)$ has coordinate set
\[
 E_i=\{\ell\in Q^\vee:\ell a_i\ne0\}.
\]
Every nonzero functional on $P$ occurs exactly $2^{q-p}$ times
among the restrictions $\ell a_i$, $\ell\in E_i$.  Consequently,
for any $i,j$ there is a coordinate bijection $E_i\to E_j$
intertwining the two marked embeddings of $P$.
We also have the exact intersection identity
\[
 U_i\cap j(Q)=ja_i(P).
 \tag{BV.4}
\]
Indeed, a vector $j(v)$ belongs to $U_i$ precisely when every
functional vanishing on $a_i(P)$ also vanishes on $v$; finite-dimensional
linear duality says exactly that $v\in a_i(P)$.

It follows from (BV.4) that the reference order induced on $U_i$
has all its vectors outside $ja_i(P)$ first, uniformly ordered,
followed by $ja_i(P)$ in the order of $P$.  The induced basis order
is independently uniform.  Restriction of a uniform order on a finite
set is uniform on each fixed subset; all fibres of the restriction
map have the same size.  The coordinate bijection just described
therefore shows that the marked ordered basis-closure type of $ja_i$
has exactly the same distribution for every $i$.  By (BV.1),
the images $a_i^*\mu$ all agree.

Lemma~\ref{A:exact-cover-averaging} now gives a rational affine
equalizer of the arbitrary family $a_i$.  Since $\mathcal A$ has
common successors, the all-family affine criterion proves rational
exactness.  The remaining properties in the statement were proved
for this same category in the earlier structural, amenability, and
parity results; no claim of integral exactness is made.
\end{proof}

\paragraph{Scope and source audit.}
The Ramsey input is Theorem~4.46 in the opened arXiv v4 of
\cite{DHubickaNesetril2019}, dated 3 September 2019:
the definitions are on printed page 75 and the theorem and complete
incidence-closure proof on printed page 76.  The padded-tuple
completion above verifies the passage to powerset orders explicitly.
The finite averaging proof uses an auxiliary exact category and
finite supports; it does not assert that its basis expansion is
precompact or invoke a general theorem about free order expansions
of vector spaces.  The universal implications in A.3(iii) and the
total-category version of C.4 remain undecided by this result.

